\documentclass[11pt]{article}
\usepackage[margin=0.8in]{geometry}
\usepackage{amsmath,amssymb,amsthm}
\usepackage{booktabs}
\usepackage{array}
\usepackage{enumitem}
\usepackage[colorlinks=true,linkcolor=blue,citecolor=blue,urlcolor=blue]{hyperref}

\numberwithin{equation}{section}
\newtheorem{theorem}{Theorem}[section]
\newtheorem{proposition}[theorem]{Proposition}
\newtheorem{lemma}[theorem]{Lemma}
\newtheorem{corollary}[theorem]{Corollary}
\newtheorem*{thm1}{Theorem 1}
\newtheorem*{thm2}{Theorem 2}
\newtheorem*{thm3}{Theorem 3}
\newtheorem*{cor4}{Corollary 4}
\newcommand{\allr}{\ref{thm:FA3}$'$}
\newtheorem*{propallr}{Proposition \ref{thm:FA3}$'$}
\theoremstyle{remark}
\newtheorem{remark}[theorem]{Remark}

\newcommand{\PP}{\mathbb{P}}
\newcommand{\CC}{\mathbb{C}}
\newcommand{\QQ}{\mathbb{Q}}
\newcommand{\ZZ}{\mathbb{Z}}
\newcommand{\NL}{\mathrm{NL}}
\newcommand{\VZ}{V_Z}
\newcommand{\cW}{\mathcal{W}}
\newcommand{\cM}{\mathcal{M}}
\newcommand{\cA}{\mathcal{A}}
\newcommand{\cX}{\mathcal{X}}
\newcommand{\Res}{\operatorname{Res}}
\newcommand{\Hess}{\operatorname{Hess}}
\newcommand{\codim}{\operatorname{codim}}
\newcommand{\T}{\mathsf{T}}
\newcommand{\lz}{\lambda_0}
\newcommand{\et}{\varepsilon_\tau}
\newcommand{\run}[1]{\medskip\noindent\textbf{#1}}
\setlist{nosep,leftmargin=2em}
\emergencystretch=1.5em

\title{\vspace{-2.5em}The special coefficient for Hodge loci of sums of two\\ complete intersections sharing equations:\\
a liaison identity and a second-order law}
\author{[author names to be added]}
\date{22 November 2026 (version 21)}

\begin{document}
\maketitle

\begin{center}
\fbox{\parbox{0.93\textwidth}{\small\textbf{Notice.} This manuscript was produced by an AI system (Claude)
during one working session, at the direction of a human user; \textbf{no human mathematician has verified
anything in it.} The working notes behind it were read by independent AI referee agents: Theorem~3 and the cases
$c=2,3$ of Corollary~4 by one (``referee~10'': no fatal error, one gap, closed in Lemma~\ref{lem:mu}); the
flat-family argument of \S\ref{sec:W} by one (``referee~11'': no fatal error, one false sentence of ours corrected).
The proof of Theorem~1 in \S\ref{sec:R} was found by referee~11 while refereeing; it and a first form of
Theorem~2 were read by one further AI referee (``referee~12'': no fatal error; one gap, closed here; it also
supplied the proof of Lemma~\ref{lem:NW}). Version~1 of this manuscript was read cold by a fourth (``referee~13'':
no fatal error; eighteen required corrections, applied here). The range-free parts of Theorem~2 and the
extension of Corollary~4 from $\tau=0$ to $\et=1$ were proposed by referee~13, checked by the coordinating AI
thread, and then read by one further AI referee in version~3 (``referee~14'': no fatal error; ten required
corrections, applied here); Theorem~2(5) was proposed by referee~14; \S\ref{sec:conj} (the formula for $A$) was found numerically by one agent, proved by the coordinating thread and
read by one AI referee (``referee~15'': no fatal error; seven repairs, applied; it also tested the intermediate
identity numerically). Version~4 was read cold by a further AI referee (``referee~16'': no fatal error; ten
required corrections, applied here; it re-derived Theorem~\ref{thm:FA}, Corollary~\ref{cor:Anz} and
Theorem~2(5)). Three statements of prior work were re-read by a second scout agent (machine digests).
Theorem~\ref{thm:FA2} was written by the coordinating thread and read by one AI referee (``referee~17'': no fatal
error; three gaps, closed by it and incorporated; it found the criterion in Corollary~\ref{cor:six});
The new parts of version~6 (\S\ref{sec:FA2} and the rewrites of version~5) were read cold by a further AI
referee (``referee~18'': no fatal error; one constant in Corollary~\ref{cor:surf} corrected; eleven corrections,
applied here; it tested the closed form of Corollary~\ref{cor:surf} numerically and proposed its second
sentence). Theorem~\ref{thm:FA3} was written by the coordinating thread and read by one AI referee
(``referee~19'': no fatal error; one gap, closed by it; it checked that the inputs do not use the range,
reproduced the exact computations behind Corollary~\ref{cor:eight}(1), (2) with its own code, and supplied
Proposition~\ref{prop:k5}). Lemma~\ref{lem:PF} (which replaces the machine-read period formula as an input) was
written by the coordinating thread and read, together with \S\ref{sec:FA3}, by one AI referee (``referee~20'': no
fatal error; it verified the analogue of the lemma for $k=0$ by hand, reproduced the computation behind
Corollary~\ref{cor:eight}(3), and asked for fifteen corrections, applied here); the cup-product formula it
rests on was looked up by a scout agent in later papers that transcribe it (machine digests; the original was
not opened). Remark~(vii) of \S\ref{sec:FA3} rests on a scout agent's digest of \cite{Kl2} and on a computation
of the coordinating thread. Theorem~\ref{thm:lin}, which collects the results for two linear subspaces, was
assembled by the coordinating thread and read, with versions~9 to~11 as a whole, by one AI referee
(``referee~21'': no fatal error; it reproduced the first-order list by exact enumeration for $d\le7$, $k\le6$ and
the second-order computations for six rows of Table~\ref{tab:lin}, gave instances that can be checked by hand,
and asked for fifteen corrections, applied here). Versions~12 to~14 (those corrections; Remark~(xi) and the
second Lean file; one further numerical run) were then read end to end by an AI referee that had no access to
the working notes or to earlier reports (``referee~22'': no false statement and no circular proof found; it
re-derived the proofs of all numbered statements except Proposition~\ref{prop:W}, ran eight independent
checks, and asked for twenty-two corrections, mostly of wording, applied here). Version~16 differs from it only in
the paragraph on Theorem~1.4 of \cite{Kl} in \S\ref{sec:intro}, rewritten after a third scout agent's digest.
Versions~17 to~21 are described in the paragraph after this notice; of them, only the changed passages of
versions~20 and~21 were read by AI referees (23, 24 and 25). Otherwise versions~15 to~21 have not been re-read. Computations are modulo primes at random instances
unless marked exact. The Lean~4 kernel certifies only the ring identities, and a few implications between polynomial equations,
named in Remark~\ref{rem:lean}, in the Checks of \S\ref{sec:FA2} and in Remark~(xi) of \S\ref{sec:lin}, and the
deduction steps of Remark~\ref{rem:lean2}, in which every geometric input is a hypothesis. On
prior work see \S\ref{sec:lit}: for linear subspaces, the consequence of Theorem~1 in the two equality cases
(the pencil $F\pm t\det M$ lies in the Hodge locus of $[\Pi_1]\mp[\Pi_2]$) is due to Kloosterman, with other
witness cycles; beyond that a finite search by an AI agent found nothing. References marked ``not opened'' were
not read by us; the others were read through machine extraction only.}}
\end{center}

{\small\noindent\textbf{Versions 17 to 21.} Version~17 adds only Remark~\ref{rem:lean2} and the mention of it in the notice. Version~18 adds only
Remarks~\ref{rem:r2} (checks by other methods, among them an implementation by an agent that had only this
paper: no discrepancy found) and~\ref{rem:pw} (a pointwise formula for the second-order term which, if correct,
removes the dependence on (CG)).
Version~19 has the final counts of that agent and the rewording of (G), (R), (PF) in \S\ref{sec:facts} which its
reading notes asked for. Version~20 moves the pointwise formula into Lemma~\ref{lem:PF} (parts (3), (4)), after an AI referee
(``referee~23'') had checked Remark~\ref{rem:pw} of version~19, and removes the hypothesis (CG) from the statements;
it also adds Proposition~\allr, which removes the hypothesis $r\ge0$ from Theorem~\ref{thm:FA3}(2), (3).
The changed passages of version~20 were then read by one more AI referee (``referee~24''): no error and no circularity found in Lemma~\ref{lem:PF}(3), (4); one gap in the first proof of
Proposition~\allr (entries of $\dot M$ of negative formal degree), closed by step (ii) of the present proof,
which is that referee's argument; nine corrections of wording, applied; and an assessment of what the numerical
checks of \eqref{eq:pw} test, now in Remark~\ref{rem:pw}. The repaired proof, the final text of Lemma~\ref{lem:PF}(3), (4) and the rewritten checks were then read by a
third AI referee (``referee~25'', who saw only version~20): no mathematical error found; one overstatement in the
checks of \eqref{eq:pw} (the test of Remark~\ref{rem:pw}(i) is blind to the shared pairs), one one-line gap in the
remark on $r<0$ and six points of wording, all corrected in version~21. Version~21 also adds the Lean files (e),
(f) of Remark~\ref{rem:lean2} and the computation of $E'$ for $r<0$ in the remark after Proposition~\allr; these
additions, and the corrections, were not re-read.\par}
\medskip

\begin{abstract}
Let $Y_1=V(f,p)$ and $Y_2=V(h,p)$ be $k$-dimensional complete intersections on a smooth hypersurface
$X=\{F=0\}\subset\PP^{2k+1}$ of degree $d$, sharing the equations $p$ and with $c$ equations each of their own, so
that $F=f^\T Mh+\sum p_ls_l$ for a $c\times c$ matrix $M$ of forms. Put $\lz=(-1)^{c-1}$ and
$\tau=\sum_{i<c}(\deg f_i+\deg h_i)-(c-1)d=d-\deg\det M$. (1) On $X$ one has
$[Y_1]+\lz[Y_2]=m''\eta^k-[R]$ ($m''$ an explicit integer), with $\eta$ the hyperplane class and $R=V(p,M^\T f,\det M)$ the Northcott residual
of $Y_1$ in the complete intersection $V(p,M^\T f)$; if $\tau\ge0$, every hypersurface $F-N\det M$,
$\deg N=\tau$, contains $R$. (2) Let $\VZ$ be the locus of hypersurfaces containing deformations of both cycles
and $\et=\dim(S/(f,h,p,s))_\tau$. For every $\lambda\ne0$ the Hodge locus of $[Y_1]+\lambda[Y_2]$ has at most
$\et$ excess tangent directions over $\VZ$ (exactly $\et$ if $(k-1)d+\tau<2k+2$; in general, for
$\lambda\ne\lz$, as many as the annihilator of $\det M$ in $S/(f,h,p,s)$ has in degree $\tau$); it equals $\VZ$ if
$\tau<0$; and for $\lambda=\lz$, $\tau\ge0$, it is smooth of codimension $\codim\VZ-\et$. (3) When $\et=1$ (for instance
$\tau=0$) and $(k-1)d+\tau<2k+2$, the second-order term of the Hodge locus of $[Y_1]+\lambda[Y_2]$ is
$\frac{\lambda-\lz}{\lambda}A$ with $A$ independent of $\lambda$; so where $A\ne0$ the locus is $\VZ$ with a double
structure for every $\lambda\notin\{0,\lz\}$. Statement (3) follows from a general principle for sums of two
classes with smooth, transversal Hodge loci and one excess direction. According to machine digests of the sources (no person has compared them): for two linear subspaces in the
two equality cases $(d,c)=(4,2),(3,3)$ the pencil $F\pm t\det M$ and the special coefficients are due to
Kloosterman, the excess count to Villaflor Loyola (at the Fermat point) and Kloosterman, and (3) is in our
note \cite{P1}; for two coplanar lines the alternative in (3) is, in the cases they treat, due to Maclean and
Dan; for disjoint curves on surfaces (2) and (3) are in our note \cite{P2}, with a different witness
curve at $\lz$. A finite literature search found, outside our own notes, neither the identity (1) in general
nor a statement of the type (3) for $c\ge2$. Finally, for $\tau=0$ in the range we prove a formula for $A$ (it uses that the constant in the period
formula for complete intersections does not depend on $X$; since version~20 this is proved in
Lemma~\ref{lem:PF} without the cup-product formula of Carlson and Griffiths): it is a nonzero constant times the derivative of the residue symbol
$\Res_S[\det M/(f,p;h,s)]$ along degree-lowering deformations of the data. It specialises, up to normalisation,
to the formulas claimed in \cite{P1} (the factor $\det C$) and \cite{P2} (where it would remove a hypothesis),
agrees modulo primes with the computed second-order term at $25$ random instances of twelve types, and shows
that $A\ne0$ for general data of every degree type with $\tau=0$ in the range that admits a smooth $X$ and
for which the second-order question is not empty. The law and the formula extend to any number $\et$ of excess directions, with $A$ replaced by a
system of quadrics on $\cA_\tau$ (Theorem~\ref{thm:FA2}); for instance, on every smooth cubic sixfold containing
two $3$-planes that meet in a line, the Hodge locus of $[\Pi_1]+\lambda[\Pi_2]$ is smooth for $\lambda=-1$ and
non-reduced for every other $\lambda\ne0$. Outside the range the law and the formula survive on the excess
directions of $\lambda\ne\lz$ (Theorem~\ref{thm:FA3}); this gives the same dichotomy for cubic eightfolds
containing two $4$-planes that meet in a plane, while for cubics of dimension $\ge10$ containing two linear
subspaces that meet in codimension two the locus is smooth for every $\lambda\ne0$ at a general point.
For two linear subspaces on a hypersurface of degree $d\ge3$ that is not a cubic surface, a cubic fourfold or a
quartic surface, the outcome is a complete list (Theorem~\ref{thm:lin}): at a general point the Hodge locus of
$[\Pi_1]+\lambda[\Pi_2]$, $\lambda\ne0,\lz$, is the locus where both subspaces deform, except in nine sporadic
cases and for coplanar lines, where it is that locus with a non-reduced
structure.
\end{abstract}

\section{Statements}\label{sec:intro}

\run{Setting.} $S=\CC[x_0,\dots,x_{2k+1}]$; $X=\{F=0\}\subset\PP^{2k+1}$ is a smooth hypersurface of degree $d$;
$\eta\in H^2(X,\ZZ)$ is the hyperplane class. Fix $c\ge1$, $k\ge c-1$, and forms
$f_0,\dots,f_{c-1}$, $h_0,\dots,h_{c-1}$, $p_c,\dots,p_k$ of degrees $d_i$, $e_j$, $\deg p_l$, all $\ge1$ and $<d$,
such that
\[
Y_1=V(f,p),\qquad Y_2=V(h,p)
\]
are complete intersections of dimension $k$ contained in $X$ and $V(f,h,p)$ has dimension $k-c$ (it is empty if
$c=k+1$). We write $I(Y_1)=(f,p)$, $I(Y_2)=(h,p)$. Then $(f,h,p)$ is a regular sequence,
$I(Y_1)\cap I(Y_2)=(p)+(f)(h)$ (Lemma~\ref{lem:alg}) and
\begin{equation}\label{eq:F}
F=f^\T Mh+\sum_l p_ls_l,\qquad g:=Mh,\quad q:=M^\T f,\qquad F=\sum_i f_ig_i+\sum_l p_ls_l=\sum_j h_jq_j+\sum_l p_ls_l,
\end{equation}
with $M$ a $c\times c$ matrix of forms, $\deg M_{ij}=d-d_i-e_j$ (and $M_{ij}=0$ if this is negative). Put
\begin{gather*}
\tau:=\sum_{i<c}(d_i+e_i)-(c-1)d=d-\deg\det M,\qquad \lz:=(-1)^{c-1},\\
\cA:=S/(f,h,p,s),\qquad \et:=\dim \cA_\tau\ \ (:=0\text{ if }\tau<0).
\end{gather*}
$\cA$ is an Artinian complete intersection (Lemma~\ref{lem:alg}) with socle degree
$\sigma=kd-2k-2+\tau$. For a rational class $\delta$ of type $(k,k)$ on $X$, $\NL(\delta)$ denotes the Hodge
locus of $\delta$: the germ at $X$, in the space of hypersurfaces of degree $d$, of the analytic space where the
flat transport of $\delta$ stays of type $(k,k)$; it is the zero scheme of the period functions $\Phi^\delta_Q$
of \S\ref{sec:facts}. We write $\delta_\lambda:=[Y_1]+\lambda[Y_2]$ for $\lambda\in\CC$ (for $\lambda\notin\QQ$, $\NL(\delta_\lambda)$ means
the zero scheme of the period functions $\Phi^{\delta_\lambda}_Q$), $V_i:=\NL([Y_i])$, $\VZ:=V_1\cap V_2$, and
identify tangent spaces with subspaces of $S_d$. Since $\VZ\subset\NL(\delta_\lambda)$ (as analytic spaces:
$\Phi^{\delta_\lambda}_Q=\Phi^{\delta_1}_Q+\lambda\Phi^{\delta_2}_Q$), the number
$\dim T\,\NL(\delta_\lambda)/T\VZ$ is called the \emph{excess} of $\lambda$. \emph{The range} is the inequality
$(k-1)d+\tau<2k+2$.

\begin{thm1}[liaison identity]
In the setting above (any $\tau$), let $R:=V(p,q_0,\dots,q_{c-1},\det M)$. Then $\tau\le d$; $R\subset X$ is
Cohen--Macaulay of pure dimension $k$ (empty iff $\tau=d$); it is linked to $Y_1$ by the complete intersection
$V(p,q)\subset X$; and
\[
[Y_1]+(-1)^{c-1}[Y_2]=m''\,\eta^k-[R]\qquad\text{in }H^{2k}(X,\ZZ)
\]
for an explicit integer $m''$ \eqref{eq:mpp}. If $\tau\ge0$, every hypersurface $\{F-N\det M=0\}$, $N\in S_\tau$,
contains $R$; hence for $N$ small it lies in $\NL(\delta_{\lz})$. Moreover $N\det M\in T V_1$ iff
$N\det M\in TV_2$ iff $N\in(f,h,p,s)_\tau$. ($R$ depends on the choice of $M$ and $s$ in \eqref{eq:F}; its class
does not.)
\end{thm1}

\begin{thm2}[first order; the coefficient $\lz$]
In the setting above:
\begin{enumerate}
\item $V_1$, $V_2$ are smooth at $X$, of codimensions $\dim(S/(f,p,g,s))_d$ and $\dim(S/(h,p,q,s))_d$. $\VZ$ is
smooth with $T\VZ=TV_1\cap TV_2$, and it is the image under $\mu(f,h,p,M,s)=f^\T Mh+\sum p_ls_l$ of a
neighbourhood of the given data. $V_1$ and $V_2$ are transversal at $X$ iff the range holds.
\item For every $\lambda\ne0$ the excess of $\lambda$ is at most $\et$; in the range it equals $\et$ for every
$\lambda\ne0$.
\item If $\tau<0$, then $\NL(\delta_\lambda)=\VZ$ for every $\lambda\ne0$.
\item If $\tau\ge0$, then $\NL(\delta_{\lz})$ is smooth of codimension $\codim\VZ-\et$, with tangent space
$T\VZ+\det M\cdot S_\tau$; near $X$ it is the set of hypersurfaces $F'-N\det M'$ with $\{F'=0\}\in\VZ$ (written as
in \eqref{eq:F} with data near the given data) and $N\in S_\tau$ small.
\item If $\tau\ge0$, then for every $\lambda\notin\{0,\lz\}$ the excess of $\lambda$ is
$\dim(0:_{\cA}\det M)_\tau=\et-\operatorname{rank}\bigl(\cdot\det M:\cA_\tau\to\cA_d\bigr)$. In particular it does
not depend on $\lambda$, and for $\tau=0$ outside the range it is $1$ if $\det M\in(f,h,p,s)$ and $0$ otherwise.
\end{enumerate}
\end{thm2}

Parts (2)--(5) need no range hypothesis; this was pointed out by referee~13, and (5) is due to referee~14.

\run{Part (3) and the literature} (machine digests; a person should confirm). Kloosterman's Theorem~1.5
(= 3.16) \cite{Kl} has, as digested, exactly our set-up and the inequality $\tau<0$ (with the standing
assumptions that $d\ge2+2/k$ and that the two primitive classes are independent), and concludes that the tangent
spaces of $\NL(\delta_\lambda)$ and of $\VZ$ have the same codimension for all but finitely many $\lambda$,
together with an equality of the loci whose quantifier we could not determine. (3) asserts the equality for
every $\lambda\ne0$, because $I_1\cap I_2=0$ in the notation of \S\ref{sec:thm2}; as digested, his own Lemma~2.11
may already give this at the level of tangent spaces. For linear subspaces and every $\lambda\ne0$ the equality
is Villaflor Loyola's theorem at the Fermat point (\cite{VL}, Theorem~3 of the arXiv version), and at a general
point he proves, as digested, the variational Hodge conjecture for these combinations of linear cycles; at
every smooth $X$ we have not found it stated. No published
example with $\tau<0$ and a larger locus for some $\lambda\ne0$ was found.

\run{Two consequences of Theorem 2(4) for cubics} (machine digests; a person should confirm before either is
called a proof of a published conjecture). For two $k$-planes on a cubic meeting in codimension one ($c=1$,
$\tau=2$): $\NL([\Pi_1]+[\Pi_2])$ is the locus of cubics containing a complete intersection of type
$(1,\dots,1,2)$ (known: \cite[Thm~5]{Mo}, \cite{Kl2}). For two $k$-planes meeting in codimension two ($c=2$,
$\tau=1$, every $k\ge2$): $\NL([\Pi_1]-[\Pi_2])$ is smooth and equal, near every smooth $X$ containing the two
planes, to the locus of cubics containing a scheme $V(p'',\text{$2\times2$ minors of a $2\times3$ matrix of linear
forms})$ with data near $(p,f,M)$: a cubic scroll for $k=2$, $\PP^1\times\PP^2\subset\PP^5$ for $k=3$, a cone over
$\PP^1\times\PP^2$ for $k\ge4$; its codimension is $\codim\VZ-\et=2\binom{k+1}3-\binom{k-1}3-(k-1)$, i.e.
$1,6,16,32,55$ for $k=2,\dots,6$. As digested, this is the object and these are the codimensions of Movasati's
conjecture on ``cubic ruled cycles'' (\cite{Mo19}, \S5, equation (12) and Table~1; \cite[Conj.~7]{Mo}), which is
stated at the Fermat point and supported there by computations to finite order; for $k=2$ the locus is the
classical divisor of cubic fourfolds containing a cubic scroll; the other coefficients are treated in
Corollary~\ref{cor:six} ($k=3$), Corollary~\ref{cor:eight} ($k=4$) and Proposition~\ref{prop:k5} ($k\ge5$). One caution, as digested by two scout agents: the introduction of \cite{Kl} states (Theorem~1.4) for
$(c,d)\in\{(2,3),(3,3),(2,4)\}$, every $\lambda\in\QQ^*$ and split $X$ that the reduced locus contains
$\NL([\Pi_1],[\Pi_2])$ properly, and for $d=3$, $k\ge5$ that the locus is reducible at $X$; read literally for
$(c,d)=(2,3)$ and $\lambda=-1$ this contradicts the smoothness asserted here. But the corresponding statement
in the body (Theorem~5.3) and its proof treat only $(d,c)=(3,3)$ and $(4,2)$, i.e. $\tau=0$ (the proof uses a
period-map argument on a K3 surface or a cubic fourfold, which is not available for $(3,2)$), and in those two
cases the special coefficient is treated separately in \S6 of \cite{Kl}. So no statement proved there, as
digested, contradicts Theorem~2(4); a person should confirm this reading. (Where we assert that the reduced
locus is $\VZ$, it is at a general $X$, not at split $X$.)

\begin{thm3}[second order, abstract]
Let $\delta_1,\delta_2$ be rational $(k,k)$-classes on a smooth hypersurface $X$ of degree $d$ in $\PP^{2k+1}$,
$\delta_\lambda=\delta_1+\lambda\delta_2$, $V_i=\NL(\delta_i)$, $\VZ=V_1\cap V_2$. Assume
\begin{enumerate}
\item[(a)] $V_1$ and $V_2$ are smooth at $X$;
\item[(b)] $T V_1+TV_2=S_d$;
\item[(c)] for every $X'\in\VZ$ near $X$ and one (equivalently every) $\lambda\ne0$,
$T_{X'}\NL(\delta_\lambda)$ contains $T_{X'}\VZ$ as a hyperplane;
\item[(d)] for one $\lz\ne0$ there is a holomorphic arc $X_t\in\NL(\delta_{\lz})$, $X_0=X$, with tangent
vector not in $T\VZ$.
\end{enumerate}
Then the space $K$ of $Q\in S_{kd-2k-2}$ whose period function $\Phi^{\delta_\lambda}_Q$ has zero differential at
$X$ is the same for all $\lambda\neq0$; $\NL(\delta_\lambda)$ lies in a smooth germ $\cM_\lambda$ of dimension
$\dim\VZ+1$ in which $\VZ=\{y=0\}$, and $\NL(\delta_\lambda)=\{y^2\upsilon_Q=0,\ Q\in K\}$; and there are a vector
$G_\lambda\in T\NL(\delta_\lambda)\smallsetminus T\VZ$ and a linear form $A$ on $K$, independent of $\lambda$,
with
\begin{equation}\label{eq:law}
\upsilon_Q(X)\,dy(G_\lambda)^2=\tfrac12\Hess\bigl(\Phi^{\delta_\lambda}_Q|_{\cM_\lambda}\bigr)(G_\lambda)
=\frac{\lambda-\lz}{\lambda}\,A(Q)\qquad(Q\in K,\ \lambda\ne0).
\end{equation}
If $A\neq0$ at $X$, then $\NL(\delta_\lambda)=\{y^2=0\}$ for every $\lambda\notin\{0,\lz\}$, and (d) can hold
for at most one coefficient. If the arcs of (d) exist through every $X'\in\VZ$ near $X$ and depend
holomorphically on $X'$, then $\NL(\delta_{\lz})=\cM_{\lz}$ is smooth of dimension $\dim\VZ+1$.
\end{thm3}

\begin{cor4}[one excess direction]
Assume $\et=1$ and the range. Then (a)--(d) hold for $\delta_i=[Y_i]$ with $\lz=(-1)^{c-1}$ and the arc
$\{F-tN\det M=0\}$, $N$ a generator of $\cA_\tau$, at every smooth $X$ of the setting. So: one excess direction
for every $\lambda\ne0$; the law \eqref{eq:law}; and $\NL(\delta_\lambda)=\VZ$ with a double structure for
$\lambda\notin\{0,\lz\}$ wherever $A\ne0$. This applies in particular in the \emph{equality case} $\tau=0$ (where
$N=1$ and necessarily $c\ge2$) when $(k-1)d<2k+2$.
\end{cor4}

With $\tau>0$, $\et=1$ holds either because $\tau=\sigma$ (exactly for $(k,d)\in\{(1,4),(2,3)\}$, where $K=0$), or
because $0<\tau<\sigma$, which forces all but one of the forms $f,h,p,s$ to be linear. With the range and $K\ne0$
only surfaces remain (referee~18: by unimodality $\cA=\CC[y]/(y^a)$ with $a\le d-1$, so
$kd-2k-2+\tau=a-1\le d-2$, which is impossible for $k\ge2$ unless $\tau=\sigma$; for $k=1$ it gives $a=d-3+\tau$
and $\tau\le2$, and $c\in\{1,2\}$ leaves the following four types): two coplanar lines ($d\ge5$); two
plane curves of degree $d-1$ in different planes, on a common surface $p=0$ of degree $d-1$ ($c=1$, $d\ge5$; the
residual lines are coplanar); a line and a disjoint plane curve of degree $d-1$ ($d\ge5$; residually, coplanar
lines again); and a line and a disjoint plane curve of degree $d-2$ ($\tau=1$, $d\ge6$; the locus at $\lz$
consists of the surfaces containing a twisted cubic, of codimension $3d-11$).

\run{Where $A\neq0$.} Theorem~\ref{thm:FA} gives a formula for $A$ when $\tau=0$, and
Corollary~\ref{cor:Anz} shows that $A\ne0$ for general data of every type with $K\ne0$; so for
$\tau=0$ the double structure in Corollary~4 is the general behaviour, not an extra hypothesis. For $\tau>0$,
$\et=1$ (the surface cases above) Corollary~\ref{cor:surf} gives $A$ explicitly: $A\ne0$ at every smooth $X$ if $d\ge6$, and for
general data if $d=5$
(for coplanar lines this is, in the cases they treat, due to Maclean and Dan). For $\et\ge2$ the number $A$ is
replaced by a system of quadrics (Theorem~\ref{thm:FA2}). Independently, modulo primes at
random instances, $A\ne0$ in every type of Table~\ref{tab:types} for which the second-order computation was run
(all rows except those with $k=4$). $A$ does vanish at special points: at the ``split'' hypersurfaces
of \cite[Def.~5.1]{Kl} (which include the Fermat hypersurface) Kloosterman's Theorem~5.3 shows, as digested,
that the reduced locus is strictly larger than $\VZ$ for every $\lambda$ in the two linear equality cases, which
by Corollary~4 would force $A=0$ there for the $k$ of the range ($d=3$: $k\le4$; $d=4$: $k\le2$); this agrees with the vanishing found in \cite{P1,P2} and would explain
why computations at the Fermat point \cite{Mo,Mo26} see no obstruction. And $K=0$, so that Corollary~4 says only
that $\NL(\delta_\lambda)$ is a smooth hypersurface in the parameter space for every $\lambda$, for skew lines on
quartic surfaces and for disjoint planes on cubic fourfolds. (The programs of \S\ref{sec:comp} compute $K$ and
$A$ in the Jacobian ring, i.e. modulo the Jacobian ideal of $F$.)

\run{The types in the range with $\tau=0$.} $(k-1)d<2k+2$ and $\tau=0$ means: $k=1$ ($c=2$, all $d\ge4$:
disjoint curves on surfaces, \cite{P2}); $k=2$ ($d=3$: two disjoint planes on a cubic fourfold, where $K=0$;
$d=4,5$: $c\le3$); $k=3$, $d=3$ ($c=3$: $f,h$ linear and one shared $p$ of degree $1$ or $2$ --- for $p$ linear,
two $3$-planes meeting in a point; $c=4$: a $3$-plane and a disjoint quadric threefold); $k=4$, $d=3$ ($c=3$;
$c=4$ with one quadric; $c=5$: two types); none for $k\ge5$. The shared equations may have any degrees in
$[1,d-1]$. For linear subspaces $\PP^k,\PP^k$ meeting in $\PP^m$ one has $c=k-m$ and $\lz=(-1)^{k-1-m}$: the cases
$(d,c)=(3,3)$ and $(4,2)$ are in \cite{P1}, where the constant in \eqref{eq:law} is computed. With $\tau>0$ and
$\et=1$: for example two coplanar lines on a surface of degree $d\ge5$ ($k=c=1$, $\tau=2$).

\run{Beyond one excess direction; two linear subspaces.} \S\ref{sec:FA2} extends the law and the formula to any
number of excess directions (Theorem~\ref{thm:FA2}) and, outside the range, to the excess directions of
$\lambda\ne\lz$ (Theorem~\ref{thm:FA3}); $A$ becomes a system of quadrics. For two $k$-dimensional linear
subspaces the outcome is Theorem~\ref{thm:lin}: when $kd-2k-2\ge1$ and $X$ is general, $\NL(\delta_\lambda)=\VZ$ as
schemes for every $\lambda\notin\{0,\lz\}$, except for the ten families of Table~\ref{tab:lin} (cubic sixfolds and
eightfolds with $c\le3$, quartic fourfolds with $c\le2$, quintic fourfolds with $c=1$, coplanar lines), where
it is $\VZ$ with a non-reduced structure; at $\lz$ (for $\tau\ge0$) the locus is smooth and larger.

\subsection{Relation to known work}\label{sec:lit}
What follows rests on machine extraction of the cited papers by AI agents (statement numbers as extracted) and
on a finite search; a person should confirm it.

\emph{Linear subspaces.} For two $k$-planes meeting in a $\PP^m$, $c=k-m$ and $\tau=2c-(c-1)d$; so $\tau<0$ iff
$c(d-2)>d$, the range of Villaflor Loyola's theorem \cite{VL} (no excess; as digested, at the Fermat point), and
$\tau=0$ iff
$(d,c)\in\{(4,2),(3,3)\}$. In these two cases Kloosterman \cite[\S6]{Kl} already considers the pencils
$F+t\det M$ (his Prop.~6.1, Cor.~6.2) and $F-t\det M$ (Prop.~6.4, Cor.~6.5) and shows that $\VZ$ has positive
codimension in the Hodge locus of $[\Pi_1]-[\Pi_2]$, resp. $[\Pi_1]+[\Pi_2]$; his witness cycles are a complete
intersection of type $(2,2,1,\dots,1)$ and a linear section of a cone over $G(2,5)$ (apparently our $W_t$ of
\S\ref{sec:W} for $c=2,3$), not $R$. That there is exactly one excess tangent direction in these two cases is
due to Villaflor Loyola at the Fermat point (appendix to \cite{Mo}, Prop.~21 and Remark~1) and, at every smooth
$X$ and for the $k$ of our range, to Kloosterman (\cite{Kl}, Example~3.17 and Theorem~1.3, as digested; Prop.~4.1
at the Fermat point); this is Theorem~2(2) for linear data. His
Theorem~1.2 gives the exceptional sets $\{0,-1\}$, $\{0,1\}$ of coefficients for large $k$, and his Theorem~1.3
leaves open, for $d=c=3$ and $k\le4$, whether the locus is larger or non-reduced; that alternative is addressed
in \cite{P1}, an AI-produced manuscript not verified by a human, and in Theorem~\ref{thm:lin} (non-reduced for
$\lambda\ne1$ at a general point). The cases $c=1$ ($\lz=1$; Dan \cite{Dan2}, Movasati \cite{Mo}, Kloosterman \cite{Kl2}) and $(d,c)=(3,2)$
($\lz=-1$, hypersurfaces containing cubic scrolls; \cite[Conj.~6, 7]{Mo}, which are conjectures) have $\tau>0$;
their special coefficients agree with $\lz=(-1)^{c-1}$ of Theorem~1. Movasati's Conjecture~1 of \cite{Mo}
(cubics, $c=3$, all coefficients) and his proposal for octic surfaces \cite{Mo26} ($k=1$, $c=2$) are the cases
analysed in \cite{P1,P2}.

\emph{Complete intersections.} Kloosterman's Theorem~1.5 \cite{Kl} treats two complete intersections sharing
equations under the strict inequality $\tau<0$. We found no treatment of $\tau\ge0$ for complete intersections,
no statement of the sign rule $\lz=(-1)^{c-1}$ in terms of $c$, and no use of liaison or of the Northcott
residual in this context. The identity $[C']=m\,\eta-[C]$ for curves linked in a surface is elementary, and
Theorem~1 is of that nature; what it adds is the identification of the right residual.

\emph{Second order.} Second-order invariants of Hodge and Noether--Lefschetz loci go back to Maclean
\cite{Mac} and are used by Duque Franco and Villaflor Loyola \cite{DFV1}. For two coplanar lines ($c=1$,
$\et=1$) the alternative ``reduced only at $\lambda=1$'' is, in the cases they treat, due to Maclean and Dan
\cite{Dan1,Dan2} (as digested); Theorem~2 reproduces Dan's codimensions $d-3$, $2d-6$, $2d-7$, and Corollary~4 applies to that case
as well. Outside our own notes we found no statement of this type in the equality case $\tau=0$ ($c\ge2$), and no
statement of the law \eqref{eq:law}.

\section{Standard facts used}\label{sec:facts}

These are cited, not re-proved, except (PF) and (NW), which are proved below; \S\ref{sec:open} lists which of
them a specialist should re-check. Up to version~19 the later statements said ``assume (PF)'', meaning the
constancy of $\kappa$, which was then derived from (CG). Since version~20 the constancy is
Lemma~\ref{lem:PF}(4), proved from (G), (A), (R) alone, and (CG) is used only for Lemma~\ref{lem:PF}(5) and for
comparisons, on which no result depends.
\begin{itemize}
\item[(G)] (Griffiths \cite{Gr}; Carlson--Griffiths \cite{CG}) With
$\Phi^\delta_Q(F')=\int_\delta\Res(Q\Omega/F'^k)$, $Q\in S_{kd-2k-2}$, the Hodge locus $\NL(\delta)$ is the zero
scheme of the $\Phi^\delta_Q$, and for $G\in S_d$:
$\Phi_Q(F+G)=\Phi_Q(F)-k\,\pi_\delta(QG)+\tfrac{k(k+1)}2\varphi_\delta(QG^2)-\dots$,
$\pi_\delta(N)=\int_\delta\Res(N\Omega/F^{k+1})$, $\varphi_\delta(N')=\int_\delta\Res(N'\Omega/F^{k+2})$. Until
\S\ref{sec:conj} only the shape (linear term $\pi_\delta(QG)$, quadratic term a nonzero multiple of
$\varphi_\delta(QG^2)$) is used; the formulas of \S\S\ref{sec:conj}, \ref{sec:FA2} and Lemma~\ref{lem:PF}(3), (4) use the
coefficients $-k$ and $\frac{k(k+1)}2$ (binomial expansion of $(F+G)^{-k}$), and the analogous expansion one pole
order higher: for a fixed $P\in S_{(k+1)d-2k-2}$ and the flat transport of $\delta$,
$\pi^{F+G}_\delta(P)=\pi_\delta(P)-(k+1)\varphi_\delta(PG)+\dots$. (Both come from differentiating under the integral
sign: near $F$ the tube of the flat transport of $\delta$ is represented by one fixed cycle in the complement of
all the hypersurfaces.) How to evaluate $\varphi_\delta$ at one
point is stated in Lemma~\ref{lem:PF}(3).
\item[(A)] The classes of the fundamental cycles of the fibres of a flat family of $k$-dimensional subschemes
of the members of a smooth family of hypersurfaces form a flat section (\cite{Fu}, \S10.1, \S19.2; not opened).
\item[(R)] Grothendieck residues for a regular sequence $\Psi$ of $2k+2$ forms: $\Res_S[\,\cdot\,/\Psi]$ is the
linear form on the socle degree of $S/(\Psi)$ which vanishes on $(\Psi)$ and takes the value $\prod_i\deg\Psi_i$ on
$\det\mathrm{Jac}(\Psi)$ (rows in the order of the sequence; so the symbol is alternating in $\Psi$); perfect
pairing; transformation law: $\Res_S[m\det A/(A\Psi)]=\Res_S[m/\Psi]$ for a square matrix $A$ of forms with
$A\Psi$ regular. Sequences written $(\phi;\psi)$ are in the order $\phi_0,\dots,\phi_k,\psi_0,\dots,\psi_k$; in
particular $\Phi^1=(f,p;g,s)$, $\Phi^2=(h,p;q,s)$ and $(f,p;h,s)$ are taken in this block order, and each symbol is
normalised on its own sequence.
\item[(SI)] The self-intersection formula for a regular embedding $Y\subset X$: $Y\cdot Y=\int_Yc_k(N_{Y/X})$
(\cite{Fu}, Cor.~6.3; not opened).
\item[(CG)] (Not used for any result since version~20; kept for Lemma~\ref{lem:PF}(5) and for the checks.
Carlson--Griffiths \cite{CG}, Theorem~3; not opened: read through two later papers that transcribe
its constant, arXiv:2402.14546 (eq.~(2.15)) and arXiv:2205.05791 (Prop.~2.1), and through Villaflor Loyola's
cup-product proposition \cite[\S6]{VL}; all machine-read by one agent) There is a constant
$c(k,d)\ne0$, depending only on $k$ and $d$, such that
$\int_X\omega_P\cup\omega_Q=c(k,d)\Res_S[PQ/(\partial_0F,\dots,\partial_{2k+1}F)]$ for all
$P,Q\in S_{(k+1)(d-2)}$, where $\omega_P:=\Res(P\Omega/F^{k+1})$.
\item[(PF)] For a complete intersection $Y=V(\phi_0,\dots,\phi_k)\subset X$ with $F=\sum_\nu\phi_\nu\psi_\nu$, all
$\phi_\nu,\psi_\nu$ of positive degree: on $S_{(k+1)d-2k-2}$,
$\pi_{[Y]}=\kappa\cdot\Res_S[\,\cdot\,/(\phi;\psi)]$ with a constant $\kappa\ne0$ which is the same for all such
data of given degrees. This is Lemma~\ref{lem:PF}(1), (2), (4) below, proved from (G), (A), (R), (SI).
Villaflor Loyola's theorem \cite[Theorem~1.2]{VL} gives, as extracted by referee~15,
$\int_Y\Res(P\Omega/F^{k+1})=\frac{(2\pi i)^k}{k!}\Res_S[P\det\mathrm{Jac}(\phi_0,\psi_0,\dots,\phi_k,\psi_k)/(\partial_0F,\dots,\partial_{2k+1}F)]$,
i.e.\ $\kappa=(-1)^{(k+1)(k+2)/2}(2\pi i)^k/k!$; this is consistent with Lemma~\ref{lem:PF} and with his value
$c(k,d)=-d(2\pi i)^{2k}/(k!)^2$, but no result of this note uses the value or the sign of $\kappa$ (the value
enters only the comparisons in the Checks of \S\ref{sec:conj}; it was verified by hand for $k=0$).
Until \S\ref{sec:conj} only $\kappa\ne0$ is used; Theorems~\ref{thm:FA}, \ref{thm:FA2}(3), \ref{thm:FA3}(2) and
Proposition~\allr use that $\kappa$ does not depend on $X$ (Lemma~\ref{lem:PF}(4)).
\item[(NW)] The Northcott--Wiebe colon formulas: Lemma~\ref{lem:NW} below, with a proof.
\item[(L)] (Peskine--Szpiro \cite{PS}) If $I$ is unmixed with Cohen--Macaulay quotient in a Gorenstein ring
(here a polynomial ring) and $J=(y):I$ for a regular sequence $(y)\subset I$ of length the height of $I$, then
the quotient by $J$ is Cohen--Macaulay of the same dimension, $I=(y):J$, and $[V(y)]=[V(I)]+[V(J)]$ as cycles (at
the generic points of $V(y)$: $\mathrm{length}(B/(0:I))=\mathrm{length}(B)-\mathrm{length}(B/I)$ in an Artinian
Gorenstein ring $B$).
\end{itemize}

\begin{lemma}[the period formula; parts (1), (2), (5) read by referee 20; parts (3), (4) added in version 20,
found by the independent implementer of Remark~\ref{rem:r2}(e); read by referee 23 in the form of Remark 10.3 of
version 19 and by referee 24 in the present form]\label{lem:PF}
Let $Y=V(\phi_0,\dots,\phi_k)\subset X$, $F=\sum_\nu\phi_\nu\psi_\nu$, $a_\nu:=\deg\phi_\nu\ge1$,
$b_\nu:=d-a_\nu\ge1$, and assume $kd-2k-2\ge0$. Put $t:=(k+1)(d-2)$ and $\varrho_Y:=\Res_S[\,\cdot\,/(\phi;\psi)]$ on $S_t$.
\begin{enumerate}
\item $\pi_{[Y]}=\kappa(X,Y)\,\varrho_Y$ on $S_t$ for a scalar $\kappa(X,Y)$, holomorphic in the data $(\phi,\psi)$.
\item $[Y]_{\mathrm{prim}}\cdot[Y]_{\mathrm{prim}}=(-1)^k\prod_\nu a_\nu b_\nu/d$. In particular
$[Y]_{\mathrm{prim}}\ne0$ and $\kappa(X,Y)\ne0$.
\item \emph{Second order at one point.} For all $P\in S_{t+d}$, with $\varphi_{[Y]}$ as in \textup{(G)}:
\begin{equation}\label{eq:pw}
(k+1)\,\varphi_{[Y]}(P)=\kappa(X,Y)\sum_{\nu=0}^{k}
\Res_S\Bigl[\frac{P}{(\phi_0,\dots,\phi_\nu^2,\dots,\phi_k;\ \psi_0,\dots,\psi_\nu^2,\dots,\psi_k)}\Bigr].
\end{equation}
\item $\kappa(X,Y)$ is the same for all data $(\phi,\psi)$ of the given degrees with $X$ smooth.
\item Assume \textup{(CG)}. Then $\kappa(X,Y)^2=-c(k,d)/d$. (Not used.)
\end{enumerate}
Parts \textup{(1)--(4)} use only \textup{(G)}, \textup{(A)}, \textup{(R)}, \textup{(SI)}.
\end{lemma}
\begin{proof}
$(\phi,\psi)$ is a regular sequence because $X$ is smooth (as in Lemma~\ref{lem:alg}), $S/(\phi,\psi)$ is Gorenstein
with socle degree $t$, and $Y\subset X$ is a regular embedding. Here $\pi_{[Y]}(P)=\int_X[Y]\cup\omega_P$ with $[Y]$
the class of the fundamental cycle, and $\omega_P\in F^kH^{2k}_{\mathrm{prim}}(X)$ (Griffiths).

(1) Every $G\in(\phi,\psi)_d$ is $\sum(\phi_\nu\dot\psi_\nu+\dot\phi_\nu\psi_\nu)$, and the hypersurfaces
$F_\epsilon=\sum(\phi_\nu+\epsilon\dot\phi_\nu)(\psi_\nu+\epsilon\dot\psi_\nu)$ contain the complete intersections
$Y_\epsilon=V(\phi+\epsilon\dot\phi)$, a flat family for small $\epsilon$. By (A), $\Phi^{[Y]}_Q(F_\epsilon)=0$ for
all $Q\in S_{kd-2k-2}$; as $F_\epsilon=F+\epsilon G+O(\epsilon^2)$, (G) gives $\pi_{[Y]}(QG)=0$. Since
$a_\nu,b_\nu<d$, $(\phi,\psi)_t=(\phi,\psi)_d\cdot S_{kd-2k-2}$, so $\pi_{[Y]}$ annihilates $(\phi,\psi)_t$; the
annihilator of $(\phi,\psi)_t$ in $S_t^*$ is the line spanned by $\varrho_Y$. The scalar $\kappa$ is holomorphic
in the data: $\pi_{[Y]}(P)$ is the integral of a holomorphic integrand over a fixed tube cycle, and $\varrho_Y(P)$
is a rational function of the coefficients of $(\phi,\psi)$ without poles where the sequence is regular, and not
identically zero in $P$. It depends only on $(X,Y)$: two presentations differ by
$\phi'=B\phi$, $\psi'=(B^{-1})^\T\psi+AB\phi$ with $A$ skew-symmetric, a transformation of determinant $1$.

(2) From $0\to N_{Y/X}\to\bigoplus\mathcal O(a_\nu)|_Y\to\mathcal O(d)|_Y\to0$,
$c_k(N_{Y/X})=\sum_{j\le k}e_j(a)(-d)^{k-j}h^k$; with (SI) and $\deg Y=\prod a_\nu=e_{k+1}(a)$:
$[Y]_{\mathrm{prim}}^2=Y\cdot Y-(\deg Y)^2/d=-\frac{\deg Y}d\sum_{j\le k+1}e_j(a)(-d)^{k+1-j}=-\frac{\deg Y}d\prod(a_\nu-d)$.
(For a line on a cubic surface: $-1-\frac13=-\frac43$; for a $k$-plane: $Y\cdot Y=(1-(1-d)^{k+1})/d$.) As it is
nonzero, $[Y]_{\mathrm{prim}}\ne0$, so $\pi_{[Y]}\ne0$ by Poincar\'e duality on primitive cohomology.

(3) For forms $\alpha_\nu,\beta_\nu$ of the degrees of $\phi_\nu,\psi_\nu$ let
$F_\epsilon=\sum_\nu(\phi_\nu+\epsilon \alpha_\nu)(\psi_\nu+\epsilon \beta_\nu)=F+\epsilon G+\epsilon^2H$, with
$G=\sum_\nu(\alpha_\nu\psi_\nu+\phi_\nu \beta_\nu)$ and $H=\sum_\nu \alpha_\nu \beta_\nu$. As in (1), $X_\epsilon=\{F_\epsilon=0\}$ is
smooth and contains the flat family $Y_\epsilon=V(\phi+\epsilon\alpha)$ for small $\epsilon$, so by (A)
$\Phi^{[Y]}_Q(F_\epsilon)=0$ for all $Q\in S_{kd-2k-2}$. By the cited fact (G), applied to the perturbation $\epsilon G+\epsilon^2H$ of $F$, the
coefficient of $\epsilon^2$ is $-k\,\pi_{[Y]}(QH)+\frac{k(k+1)}2\varphi_{[Y]}(QG^2)$; hence by (1)
\[
(k+1)\,\varphi_{[Y]}(QG^2)=2\,\pi_{[Y]}(QH)=2\kappa\,\Res_S[QH/(\phi;\psi)].
\]
(The right side depends on $(\alpha,\beta)$ only through $G$: two choices with the same $G$ differ by a Koszul
syzygy of $(\phi,\psi)$, which changes $H$ by an element of $(\phi,\psi)$.)
Let $\Phi_{(\nu)}$ be the sequence $(\phi;\psi)$ with $\phi_\nu$, $\psi_\nu$ replaced by their squares; it is
regular, with socle degree $t+d$, and it is obtained from $(\phi;\psi)$ by a diagonal matrix of determinant
$\phi_\nu\psi_\nu$. Modulo $\Phi_{(\nu)}$ we have $G\equiv \alpha_\nu\psi_\nu+\phi_\nu \beta_\nu$, hence
$G^2\equiv2\alpha_\nu \beta_\nu\phi_\nu\psi_\nu$, and by (R)
$\Res_S[QG^2/\Phi_{(\nu)}]=2\Res_S[Q\alpha_\nu \beta_\nu/(\phi;\psi)]$. Summing over $\nu$: \eqref{eq:pw} holds for
$P=QG^2$. Both sides are then quadratic forms in $G\in(\phi,\psi)_d$ which agree, so their polarisations agree:
\eqref{eq:pw} holds for $P=QGG'$ with $G,G'\in(\phi,\psi)_d$, $Q\in S_{kd-2k-2}$. These products span
$[(\phi,\psi)^2]_{t+d}$ (a generator $\xi\xi'm$, with $\xi,\xi'$ members of $(\phi;\psi)$ and $m$ a monomial, is of
this form, since $m$ factors into monomials of degrees $d-\deg\xi$, $d-\deg\xi'$, $kd-2k-2$), and
$[(\phi,\psi)^2]_{t+d}=S_{t+d}$: with $I=(\phi,\psi)$, $S/I$ vanishes above degree $t$, so $I_{t+d}=S_{t+d}$; and
$I/I^2\cong\bigoplus_\xi(S/I)(-\deg\xi)$ ($I$ is generated by a regular sequence), which vanishes in degrees
above $t+\max_\xi\deg\xi\le t+d-1$, so $(I^2)_{t+d}=I_{t+d}$.

(4) Along the same family differentiate $\pi^{X_\epsilon}_{[Y_\epsilon]}(P)=\kappa(\epsilon)\Res_S[P/(\phi+\epsilon\alpha;\psi+\epsilon\beta)]$,
$P\in S_t$, at $\epsilon=0$ ($\kappa(\epsilon):=\kappa(X_\epsilon,Y_\epsilon)$ is holomorphic by (1)). On the left,
by (G) one order higher (the derivative of the period of a fixed form over a flat class in the direction $G$), we
get $-(k+1)\varphi_{[Y]}(PG)$. On the right, by the rule
$\frac{d}{d\epsilon}\Res_S[P/(\dots,u+\epsilon\dot u,\dots)]=-\Res_S[P\dot u/(\dots,u^2,\dots)]$ for a member $u$ of a
regular sequence (which follows from (R) and $(u+\epsilon\dot u)(u-\epsilon\dot u)=u^2-\epsilon^2\dot u^2$; the
argument is written out in the proof of Theorem~\ref{thm:FA}(1)) and (R) again,
\[
\dot\kappa\,\Res_S\Bigl[\frac{P}{(\phi;\psi)}\Bigr]-\kappa\sum_\nu\Res_S\Bigl[\frac{PG}{\Phi_{(\nu)}}\Bigr]
\]
(for the pair $\nu$: $PG\equiv P(\alpha_\nu\psi_\nu+\phi_\nu \beta_\nu)$ modulo the other members, and
$\Res_S[P\alpha_\nu\psi_\nu/\Phi_{(\nu)}]=\Res_S[P\alpha_\nu/(\dots,\phi_\nu^2,\dots;\psi)]$, likewise for $\beta_\nu$). By
\eqref{eq:pw} for $PG\in S_{t+d}$ the two expressions differ by $\dot\kappa\Res_S[P/(\phi;\psi)]$, so
$\dot\kappa=0$. As $(\alpha,\beta)$ is an arbitrary tangent vector at an arbitrary point of the space of data, which is
Zariski-open in an affine space, hence connected, $\kappa$ is constant on it.

(5) $(\partial_iF)_i=\mathrm{Jac}(\phi_0,\psi_0,\dots,\phi_k,\psi_k)^\T(\psi_0,\phi_0,\dots,\psi_k,\phi_k)^\T$, so by (R),
with $J_Y:=\det\mathrm{Jac}(\phi_0,\psi_0,\dots)$, $\rho_J:=\Res_S[\,\cdot\,/(\partial F)]$ and
$\varsigma:=(-1)^{(k+1)(k+2)/2}$: $\rho_J(PJ_Y)=\varsigma\,\varrho_Y(P)$, and $\rho_J(J_Y^2)=(-1)^{k+1}\prod a_\nu b_\nu$ (only the
$k+1$ transpositions $\phi_\nu\leftrightarrow\psi_\nu$ matter; $\Res_S[\det\mathrm{Jac}(g)/(g)]=\prod\deg g_i$). Write
$[Y]_{\mathrm{prim}}\equiv\omega_b$ modulo $F^{k+1}$ with $b\in S_t$; then
$\pi_{[Y]}(P)=\int_X\omega_b\cup\omega_P=c(k,d)\,\rho_J(bP)$ by (CG) (note $F^k\cup F^{k+1}=0$), and
$\pi_{[Y]}(P)=\kappa\varsigma\,\rho_J(J_YP)$ by (1); so $c(k,d)\,b\equiv\kappa\varsigma J_Y$ modulo the Jacobian
ideal, and
$[Y]_{\mathrm{prim}}^2=c(k,d)\,\rho_J(b^2)=\frac{\kappa^2}{c(k,d)}\rho_J(J_Y^2)=\frac{\kappa^2}{c(k,d)}(-1)^{k+1}\prod a_\nu b_\nu$.
Compare with (2).
\end{proof}

For the two cycles $Y_1,Y_2$ of the setting the constants are \emph{equal} if $\tau\ge0$: see Lemma~\ref{lem:det}. Referee~20
verified the analogue of the lemma for $k=0$ (excluded by its hypothesis) by hand ($d$ points on a line; $\kappa=-1$, $c(0,d)=-d$; part (1) holds
there by inspection), part (2) against nine classical self-intersection numbers, and the residue identities
of (5) exactly for curves on quartic surfaces. Parts (3), (4) were read by referees~23 and~24 (Remark~\ref{rem:pw}).

\section{Algebra and first order}\label{sec:first}

\begin{lemma}\label{lem:alg}
In the setting: (i) $I(Y_1)\cap I(Y_2)=(p)+(f)(h)$, whence \eqref{eq:F}. (ii) The sequences
$\Phi^1:=(f,p;g,s)$, $\Phi^2:=(h,p;q,s)$ and $(f,h,p,s)$ consist of $2k+2$ forms of positive degree and are
regular. (iii) $\deg\det M=cd-\sum_{i<c}(d_i+e_i)=d-\tau$. (iv) The socle degree of
$S/\Phi^i$ is $(k+1)d-2k-2$ and that of $\cA=S/(f,h,p,s)$ is $\sigma=kd-2k-2+\tau$.
\end{lemma}
\begin{proof}
(i) Modulo $(p)$, $f$ and $h$ are regular sequences with $(f,h)$ regular, so
$\mathrm{Tor}_1(S/(p,f),S/(p,h))=0$ over $S/(p)$, i.e. $(f)\cap(h)=(f)(h)$ there. (ii) At a common zero of one
of these sequences $\nabla F=0$, because $F$ is a sum of products of two members of the sequence (resp.
$F\in(f)(h)+(p)(s)$); $X$ is smooth, so there is none. So none of the $2k+2$ forms is zero (otherwise the
remaining $2k+1$ would have a common zero), all have positive degree ($1\le d_i,e_j,\deg p_l<d$), and $2k+2$
forms of positive degree in $2k+2$ variables without common zero form a regular sequence. (iii) Each term of
$\det M$ has degree $\sum_i(d-d_i-e_{\pi(i)})$. (iv) Sum of the degrees minus $2k+2$.
\end{proof}

\begin{lemma}\label{lem:T}
(i) $TV_1=(\Phi^1)_d$, $TV_2=(\Phi^2)_d$, and $V_1,V_2$ are smooth at $X$: $V_1$ is the image of a neighbourhood
of $(f,p,g,s)$ under $\mu_1(f,p,g,s)=\sum f_ig_i+\sum p_ls_l$. (ii) $TV_1+TV_2=(f,h,p,s)_d$; this is $S_d$ iff
$d>\sigma$, i.e. iff the range holds.
\end{lemma}
\begin{proof}
(i) If $kd-2k-2<0$ there are no period functions and $(\Phi^1)_d=S_d$. Otherwise, by Lemma~\ref{lem:PF}(1),
(2), $\pi_1:=\pi_{[Y_1]}=\kappa_1\Res_S[\,\cdot\,/\Phi^1]$ with $\kappa_1\ne0$, so by (G) and (R) the Zariski tangent space of
$V_1$ is $\{G:\pi_1(QG)=0\ \forall Q\}=(\Phi^1)_d$ (perfect pairing of $S/\Phi^1$ between the degrees $d$ and
$kd-2k-2$). The differential of $\mu_1$ has image
$(f)_d+(g)_d+(p)_d+(s)_d=(\Phi^1)_d$, of codimension $\dim(S/\Phi^1)_d$, a number depending only on the degrees
as long as $X$ is smooth; so $\mu_1$ has constant rank and the image of a small neighbourhood is a smooth germ
with tangent space $(\Phi^1)_d$. It lies in $V_1$ by (A) (the hypersurface $\mu_1(f',p',g',s')$ contains the
deformation $V(f',p')$ of $Y_1$), and a germ of analytic space containing a smooth germ of dimension equal to
its embedding dimension coincides with it. (ii) $g\in(h)$, $q\in(f)$, and (iv) of Lemma~\ref{lem:alg}.
\end{proof}

\begin{lemma}[referee 10; range-free form after referee 13]\label{lem:mu}
$(\Phi^1)\cap(\Phi^2)=(p,s)+(g)+(q)+(f)(h)$. Consequently $\VZ=V_1\cap V_2$ is smooth with
$T\VZ=TV_1\cap TV_2$, and the map $\mu(f,h,p,M,s)=f^\T Mh+\sum p_ls_l$ is a submersion onto $\VZ$ near
$(f,h,p,M,s)$: every $X'\in\VZ$ near $X$ is of the form \eqref{eq:F} with data depending holomorphically on
$X'$, and the setting holds at $X'$.
\end{lemma}
\begin{proof}
$\supseteq$ is clear. Let $G=\sum a_if_i+\sum b_ig_i\in(f,g)\cap(h,q)$ in $S/(p,s)$. Then
$G-\sum b_ig_i\in(f)\cap((h)+(q))=(q)+((f)\cap(h))$ by the modular law ($(q)\subset(f)$), and
$(f)\cap(h)=(f)(h)$ modulo $(p,s)$ because $(p,s,f,h)$ is regular. The image of $d\mu$ is
$[(g)+(q)+(s)+(p)+(f)(h)]_d=TV_1\cap TV_2$, of dimension depending only on the degrees; so $\mu$ has constant
rank and its local image is a smooth germ of dimension $\dim(TV_1\cap TV_2)$. It lies in the analytic space
$V_1\cap V_2$ (the $V_i$ are smooth, hence reduced, and contain the image), whose Zariski tangent space is
$TV_1\cap TV_2$; and a germ of analytic space containing a smooth germ of dimension equal to its embedding
dimension coincides with it.
Smoothness of $X'$ and $\dim V(f',h',p')=k-c$ are open conditions.
\end{proof}

\begin{lemma}[Northcott \cite{No}, Wiebe \cite{Wi}; proof by referee 12]\label{lem:NW}
Let $x=(x_1,\dots,x_c)$ and $y=(y_1,\dots,y_c)$ be regular sequences in a commutative ring $B$, $(x)\ne B$,
$y=x\cdot M$ for a $c\times c$ matrix $M$ over $B$, $\Delta=\det M$. Then multiplication by $\Delta$ is an
injective map $B/(x)\to B/(y)$ with image $((y):(x))/(y)$. Equivalently $(y):(x)=(y)+(\Delta)$ and
$(y):\Delta=(x)$; hence $(y):((y)+(\Delta))=(x)$ and $\Delta\notin(y)$.
\end{lemma}
\begin{proof}
$\Lambda^\bullet M$ is a morphism of Koszul complexes $K(y)\to K(x)$ lifting the surjection
$\pi:B/(y)\to B/(x)$, and both complexes are free resolutions. Applying $\mathrm{Hom}(-,B)$ and taking $H^c$,
the map $\pi^*:\mathrm{Ext}^c(B/(x),B)\to\mathrm{Ext}^c(B/(y),B)$ is, in the Koszul bases, $B/(x)\to B/(y)$,
$1\mapsto\Delta$. On the other hand $\mathrm{Ext}^c_B(N,B)\cong\mathrm{Hom}_B(N,B/(y))$ naturally in the
$B/(y)$-module $N$ (\cite[Lemma~1.2.4]{BH}), and under this isomorphism $\pi^*$ is
$\mathrm{Hom}(B/(x),B/(y))\to B/(y)$, $\phi\mapsto\phi(1)$: injective with image $((y):(x))/(y)$. The two
descriptions differ by units.
\end{proof}
The classical statement is \cite[Theorem~2.5]{CPV}, \cite[Prop.~1.2]{SS} (as extracted).

\begin{lemma}\label{lem:det}
For $m\in S_\sigma$: $\Res_S[m\det M/\Phi^1]=\Res_S[m/(f,p;h,s)]=:\varrho(m)$ and
$\Res_S[m\det M/\Phi^2]=(-1)^c\varrho(m)$. If $\tau\ge0$ and $kd-2k-2\ge0$, the constants $\kappa(X,Y_1)$,
$\kappa(X,Y_2)$ of Lemma~\ref{lem:PF} are equal; call them $\kappa$. Hence for $N\in S_\tau$ and
$Q\in S_{kd-2k-2}$: $\pi_1(QN\det M)=\kappa\,\varrho(QN)$ and $\pi_2(QN\det M)=(-1)^c\kappa\,\varrho(QN)=-\lz\kappa\,\varrho(QN)$;
and $N\det M\in TV_i$ iff $N\in(f,h,p,s)_\tau$ $(i=1,2)$. In particular $\det M\not\equiv0$, and
$\det M\notin TV_i$ when $\tau=0$.
\end{lemma}
\begin{proof}
$\Phi^1=(f,p;Mh,s)$ is obtained from $(f,p;h,s)$ by the matrix $\mathrm{diag}(1,1,M,1)$ of determinant
$\det M$; transformation law (R). $\Phi^2=(h,p;M^\T f,s)$ is obtained in the same way from $(h,p;f,s)$, which
differs from $(f,p;h,s)$ by $c$ transpositions. The pairing $\cA_\tau\times\cA_{\sigma-\tau}\to\CC$ induced by
$\varrho$ is perfect and $\sigma-\tau=kd-2k-2$. Equality of the constants $\kappa_i:=\kappa(X,Y_i)$: the first
assertions of Theorem~1 (proved in \S\ref{sec:R} without this lemma) and (G) give
$(u_1+\lz u_2)(N\det M)=0$, i.e. $(\kappa_1+\lz(-1)^c\kappa_2)\varrho(QN)=(\kappa_1-\kappa_2)\varrho(QN)=0$ for all
$Q\in S_{kd-2k-2}$, $N\in S_\tau$, and $\varrho\ne0$ on $S_{kd-2k-2}S_\tau=S_\sigma$ (for $\tau<0$ the equality is
neither proved nor used).
\end{proof}

\section{The residual \texorpdfstring{$R$}{R}: proof of Theorem 1}\label{sec:R}

\emph{Step 1.} $F=\sum_jh_jq_j+\sum_lp_ls_l\in(p,q)$, so $V(p,q)\subset X$ and $R\subset X$; and
$F-N\det M\in I(R)$ for every $N\in S_\tau$ when $\tau\ge0$.

\emph{Step 2.} In the graded ring $S/(p)$, $f$ and $q=M^\T f$ are regular sequences ($(p,f)$ cuts out $Y_1$;
$(p,q)$ is part of $\Phi^2$: this is where the smoothness of $X$ enters). By Lemma~\ref{lem:NW},
$(p,q):(p,f)=(p,q,\det M)=I(R)$ and $(p,q):\det M=(p,f)$; in particular $\det M\notin(p,q)$, so $\det M\not\equiv0$
and $\tau\le d$. If $\deg\det M=0$ (i.e. $\tau=d$), then $R=\emptyset$ and $V(p,q)=Y_1$. Otherwise $I(R)$ is
unmixed of height $k+1$, $R$ is Cohen--Macaulay of pure dimension $k$ by (L), and at every minimal prime of
$(p,q)$ the lengths of the local rings of $Y_1$ and $R$ add up to that of $V(p,q)$. So
\begin{equation}\label{eq:link}
[V(p,q)]=[Y_1]+[R]
\end{equation}
as cycles with multiplicities ($Y_1$ and $R$ may be non-reduced and may share components).

\emph{Step 3.} For $i=0,\dots,c$ let $\Theta_i:=V(p,h_0,\dots,h_{i-1},q_i,\dots,q_{c-1})$: a complete intersection of
dimension $k$ ($k+1$ distinct members of $\Phi^2$), $\Theta_0=V(p,q)$, $\Theta_c=Y_2$. For $i<c$ the $k$ forms
$(p,h_0,\dots,h_{i-1},q_{i+1},\dots,q_{c-1})$ cut $X$ in a complete intersection $C_i$ of dimension $k$, on which
$F\equiv h_iq_i$; so $C_i=V(\dots,h_iq_i)$ has cycle $[\Theta_{i+1}]+[\Theta_i]$ (additivity of length for a product
of two non-zero-divisors), and $[C_i]=m_i\eta^k$ with
$m_i=\prod_l\deg p_l\cdot\prod_{j<i}e_j\cdot\prod_{j>i}(d-e_j)$. Hence
\begin{equation}\label{eq:mpp}
[V(p,q)]=(-1)^c[Y_2]+m''\eta^k,\qquad m''=\sum_{i<c}(-1)^im_i .
\end{equation}

\emph{Step 4.} \eqref{eq:link} and \eqref{eq:mpp} give $[Y_1]+(-1)^{c-1}[Y_2]=m''\eta^k-[R]$.

\emph{Step 5.} The hypersurfaces $X_N=\{F-N\det M=0\}$ ($N\in S_\tau$) all contain $R$. Let $U$ be a convex
neighbourhood of $F$ in $S_d$ consisting of smooth hypersurfaces. For $X_N\in U$ the segment from $X$ to $X_N$
is a family $\cX\to[0,1]$ of smooth hypersurfaces $X_t$ containing the fixed scheme $R$; for each component
$R_i$ of $R$ the inclusions $R_i\hookrightarrow X_t\hookrightarrow\cX$ for different $t$ are homotopic and
$X_t\hookrightarrow\cX$ is a homotopy equivalence, so the homology class of $R_i$, and with the fixed
multiplicities that of $R$, is a flat section. Hence the transport of $\delta_{\lz}=m''\eta^k-[R]$ to $X_N$ is
$m''\eta^k-[R]$, a Hodge class, and $X_N\in\NL(\delta_{\lz})$. The last assertion of Theorem~1 is
Lemma~\ref{lem:det}. \qed

\begin{remark}\label{rem:lean}
For a generic $4\times4$ matrix ($c=4$, no $p$) the inclusion $\det M\cdot(f)\subset(q)$ (Cramer), the identity
$f^\T Mh=\sum h_jq_j$, and the identities of Proposition~\ref{prop:W}(i) are checked by the Lean~4 kernel as
ring identities (file \texttt{GenFamilyCert.lean}, 25 theorems). For $c=2,3$ the analogous identities are in
\texttt{OcticDetCert.lean} and \texttt{PfaffianCert.lean}, in the notation of \cite{P2}, \cite{P1}. The converse
inclusion in Lemma~\ref{lem:NW} is not machine-checked.
\end{remark}

\begin{remark}
For $c=1$: $M=(M_{00})$, $R=V(p,M_{00})$, and the identity reads $[Y_1]+[Y_2]+[R]=m''\eta^k$, the decomposition
of $X\cap V(p)=V(p,f_0M_{00}h_0)$. For two skew lines on a quartic surface ($c=2$) $R$ is a curve of degree $8$
and genus $7$; for two disjoint planes on a cubic fourfold ($c=3$) a surface of degree $7$.
\end{remark}

\section{Proof of Theorem 2}\label{sec:thm2}

(1) is Lemmas~\ref{lem:T} and~\ref{lem:mu}.

(2) Let $u_i(G):=\pi_i(\,\cdot\,G)\in S_{kd-2k-2}^*$ and $I_i:=u_i(S_d)=((\Phi^i)_{kd-2k-2})^\perp$; then
$TV_i=\ker u_i$ and $T\NL(\delta_\lambda)=\ker(u_1+\lambda u_2)$. If $G\in\ker(u_1+\lambda u_2)$, then
$u_1(G)=-\lambda u_2(G)\in I_1\cap I_2$, and $u_1(G)=0$ forces $u_2(G)=0$; so $G\mapsto u_1(G)$ embeds
$T\NL(\delta_\lambda)/T\VZ$ into $I_1\cap I_2=((f,h,p,s)_{kd-2k-2})^\perp\cong(\cA_{\sigma-\tau})^*$, of dimension
$\et$ by the symmetry of the Hilbert function of the Gorenstein ring $\cA$ (and $0$ if $\tau<0$, since then
$kd-2k-2>\sigma$). In the range, $(u_1,u_2):S_d\to I_1\oplus I_2$ is onto (Lemma~\ref{lem:T}(ii)), so the
embedding is onto $I_1\cap I_2$.

(3) If $\tau<0$, then by (2) $T\NL(\delta_\lambda)=T\VZ$; as $\NL(\delta_\lambda)$ contains the smooth germ
$\VZ$, the two coincide.

(4) By Theorem~1 the vectors $N\det M$ are tangent to arcs inside $\NL(\delta_{\lz})$, so they lie in its Zariski
tangent space, and modulo $T\VZ$ they span a space of dimension $\et$ (Lemma~\ref{lem:det}); with (2),
$T\NL(\delta_{\lz})=T\VZ+\det M\cdot S_\tau$. Choose a complement $S'$ of $(f,h,p,s)_\tau$ in $S_\tau$. By
Lemma~\ref{lem:mu} the map $\Psi(X',N)=F'-N\det M'$ ($X'\in\VZ$ near $X$, $N\in S'$ small) is holomorphic with
injective differential at $(X,0)$. Its image lies in $\NL(\delta_{\lz})$: $Y'_1=V(f',p')$ and $Y'_2=V(h',p')$ are
flat deformations of $Y_1,Y_2$, so by (A) the transport of $\delta_{\lz}$ from $X$ to $X'$ is
$[Y'_1]+\lz[Y'_2]$, which by Theorem~1 at $X'$ equals $m''\eta^k-[R']$, $R'=V(p',M'^\T f',\det M')$ ($m''$ depends
only on the degrees); and the segment $F'-tN\det M'$, $0\le t\le1$, stays in the convex neighbourhood $U$ of
Step~5 of \S\ref{sec:R} and all its members contain $R'$. No flatness of the family of the $R'$ is used. So
$\NL(\delta_{\lz})$ contains a smooth germ of dimension equal to its embedding dimension, and equals it. \qed

(5) (Referee~14.) Let $\lambda\notin\{0,\lz\}$ and $\gamma:=1/\lz-1/\lambda\ne0$. By Lemma~\ref{lem:det},
$N\mapsto u_1(N\det M)$ maps $S_\tau$ onto $I_1\cap I_2$ with kernel $(f,h,p,s)_\tau$, and by Theorem~1,
$u_2(N\det M)=-u_1(N\det M)/\lz$. Let $G\in\ker(u_1+\lambda u_2)$ and choose $N$ with $u_1(N\det M)=u_1(G)$. Then
$G':=G-N\det M$ has $u_1(G')=0$ and $u_2(G')=\gamma u_1(G)$, so $G'\in TV_1$ and $G'+\gamma\lz N\det M\in TV_2$; hence
$N\det M\in TV_1+TV_2=(f,h,p,s)_d$. Conversely, if $N\det M=H_1+H_2$ with $H_i\in TV_i$, then
$G:=N\det M-\gamma\lz H_1$ lies in $\ker(u_1+\lambda u_2)$ and $u_1(G)=u_1(N\det M)$. So the embedding of (2)
identifies $T\NL(\delta_\lambda)/T\VZ$ with $((f,h,p,s):\det M)_\tau/(f,h,p,s)_\tau=(0:_\cA\det M)_\tau$. \qed

\begin{remark}
(i) Lemma~\ref{lem:NW} applied in $S$ to $x=(f,p,h,s)$, $y=\Phi^i$ gives $(\Phi^i):\det M=(f,h,p,s)$, a
residue-free proof of the last assertion of Theorem~1. (ii) By Theorem~1, $\NL(\delta_{\lz})=\NL([R])$ near $X$,
and for $\tau\ge0$ Theorem~2(4) says that this is the locus of hypersurfaces containing a scheme
$V(p'',M''^\T f'',\det M'')$ of the same determinantal form as $R$, with data near $(p,f,M)$ (conversely, a form
of $I(R'')_d$ near $F$ is $\sum p''s''+\sum a_jq''_j-N\det M''$ with $(a,s'',N)$ near $(h,s,0)$, since $\dim I(R'')_d$
is constant by linkage): the Hodge locus of the residual is as small as it can be. (iii) Theorem~2(5) agrees with the instances outside the range that we
computed (\S\ref{sec:comp} and the Checks of \S\ref{sec:FA3}: excess $\et$ at $\lz$ and $\et-1$ at the other tested values) and, as far as the
digests allow a comparison, with the thresholds in Kloosterman's theorems: for $(d,c)=(3,2)$ multiplication by
$\det M$ from $\cA_1$ to $\cA_3$ is injective for general data if and only if $k\ge5$ (if: Proposition~\ref{prop:k5}; only if:
$\dim\cA_1=k-1>\binom{k-1}3=\dim\cA_3$ for $2\le k\le4$); for $(d,c)=(3,3)$, $\cA_3\ne0$ iff $k\ge5$; and at
hypersurfaces with all entries of $M$ in $(f,h,p,s)$ the excess is $\et$ for every $\lambda$ and every $k$.
\end{remark}

\section{The second-order principle: proofs of Theorem 3 and Corollary 4}\label{sec:G}

Notation as in Theorem~3: $\pi_i,\varphi_i$ for $\delta_i$; $u_i(G)=\pi_i(\,\cdot\,G)$; $I_i=u_i(S_d)$;
$T_i=TV_i=\ker u_i$.

\emph{Step 1.} By (b), $(u_1,u_2):S_d\to I_1\oplus I_2$ is onto with kernel $T\VZ$, and
$T\NL(\delta_\lambda)/T\VZ\cong I_1\cap I_2$ for every $\lambda\ne0$. (a) and (b) are open conditions, so the same
holds at every $X'\in\VZ$ near $X$: $\dim T_{X'}\NL(\delta_\lambda)/T_{X'}\VZ=\dim(I_1\cap I_2)(X')$ for every
$\lambda\ne0$; this is the equivalence asserted in (c). By (c), $I_1\cap I_2$ is a line $\CC\omega$. Since
$(u_1,u_2)$ is onto, the image of $u_1+\lambda u_2$ is $I_1+I_2$; hence $K=(I_1+I_2)^\perp$ is independent of
$\lambda$: for $Q\in K$ the functions $\Phi^{\delta_1}_Q$ and $\Phi^{\delta_2}_Q$ both have zero differential at
$X$.

\emph{Step 2.} Choose functions among the $\Phi^{\delta_\lambda}_Q$ with independent differentials cutting out a
smooth germ $\cM_\lambda\supset\NL(\delta_\lambda)$ of dimension $\dim\VZ+1$. For $Q\in K$,
$\phi_Q:=\Phi^{\delta_\lambda}_Q|_{\cM_\lambda}$ vanishes on $\VZ=\{y=0\}$ and, by (c) at the points of $\VZ$
(where $T\NL(\delta_\lambda)=T\cM_\lambda$), has zero differential along $\VZ$: $\phi_Q=y^2\upsilon_Q$. As
$d\Phi^{\delta_\lambda}_Q(X)=0$, $\tfrac12\Hess\phi_Q(G)=c_0\,\varphi_{\delta_\lambda}(QG^2)$ for
$G\in T\NL(\delta_\lambda)$, $c_0=k(k+1)/2$.

\emph{Step 3 (key lemma).} For $Q\in K$ and $G\in T_1$: $\varphi_1(QG^2)=0$; likewise for $T_2$, $\varphi_2$.
Indeed $\Phi^{\delta_1}_Q$ vanishes on the smooth germ $V_1$ and has zero differential at $X$ on all of $S_d$
(Step~1); take an arc $\gamma$ in $V_1$ with $\gamma'(0)=G$: $0=(\Phi^{\delta_1}_Q\circ\gamma)''(0)=
\Hess\Phi^{\delta_1}_Q(G,G)+d\Phi^{\delta_1}_Q(\gamma''(0))=2c_0\varphi_1(QG^2)$.

\emph{Step 4.} Choose $G_1\in T_2$ with $u_1(G_1)=\omega$ and $G_2\in T_1$ with $u_2(G_2)=\omega$ (possible by
(b)). Then $G_\lambda:=G_1-G_2/\lambda\in T\NL(\delta_\lambda)\smallsetminus T\VZ$, and for $Q\in K$, by Step~3,
\[
\varphi_{\delta_\lambda}(QG_\lambda^2)=\bigl[\varphi_1(QG_1^2)-2\varphi_2(QG_1G_2)\bigr]
+\frac1\lambda\bigl[\varphi_2(QG_2^2)-2\varphi_1(QG_1G_2)\bigr]=:A_0(Q)+\frac{B_0(Q)}{\lambda}.
\]

\emph{Step 5.} Along the arc of (d), $\phi_Q=y^2\upsilon_Q$ vanishes and $y$ has a simple zero; so $\upsilon_Q(X)=0$ at
$\lambda=\lz$ for all $Q\in K$, i.e. $A_0+B_0/\lz=0$, and
$\varphi_{\delta_\lambda}(QG_\lambda^2)=\frac{\lambda-\lz}{\lambda}A_0(Q)$. Put $A:=c_0A_0$. The remaining
assertions: if $A\ne0$, some $\upsilon_Q$ is a unit for $\lambda\notin\{0,\lz\}$; two coefficients satisfying (d) would
force $A_0=B_0=0$; and if the arcs exist through all $X'\in\VZ$, the map $(X',t)\mapsto X'_t$ has injective
differential and image in $\NL(\delta_{\lz})\subset\cM_{\lz}$, of the same dimension. \qed

\begin{proof}[Proof of Corollary 4]
(a), (b): Lemma~\ref{lem:T}, in the range. (c): Theorem~2(2) with $\et=1$, at every $X'\in\VZ$ near $X$ by
Lemma~\ref{lem:mu} ($\et$ depends only on the degrees). (d): Theorem~1 with $N$ a generator of $\cA_\tau$, so
that $N\det M\notin T V_1\supset T\VZ$ (Lemma~\ref{lem:det}); through every $X'$ by Lemma~\ref{lem:mu}. If
$\tau=0$ then $c\ge2$, since for $c=1$ one has $\tau=d_0+e_0\ge2$.
\end{proof}

\section{A second description at \texorpdfstring{$\lz$}{lambda0}: the family \texorpdfstring{$W_t$}{Wt}}\label{sec:W}

Assume $\tau=0$, $c\ge2$, $n:=c-1$. Let $a_j$ be the cofactor of $M_{jn}$ in $M$,
$w_j(t):=f_jh_n-t\,a_j$ (homogeneous of degree $d_j+e_n$), $I(t):=(p,w_0,\dots,w_n,q_0,\dots,q_{n-1})$,
$\Sigma:=V(p,f,h_n)$ (of dimension $k-1$), $\Theta:=V(p,h_n,q_0,\dots,q_{n-1})$.

\begin{proposition}[sketch; details in the working notes, read by referee 11]\label{prop:W}
(i) $F-t\det M=\sum_{l<n}h_lq_l+\sum_jM_{jn}w_j+\sum_l p_ls_l\in I(t)$; $h_nq_l=\sum_jM_{jl}w_j$ for $l<n$; and
$a_if_j-a_jf_i\in(q_0,\dots,q_{n-1})$. (ii) $I(0)=I(Y_1)\cap I(\Theta)$, without embedded components, with cycle
$[Y_1]+[\Theta]$, and $[\Theta]=(-1)^{c-1}[Y_2]+m'\eta^k$ for an integer $m'$. (iii) Let $I\subset S[t]$ be the ideal
generated by $I(t)$ and $\cW\subset\PP^{2k+1}\times\mathbb{A}^1$ the
scheme-theoretic closure of $V(I)\smallsetminus(\Sigma\times\mathbb{A}^1)$ over $t\neq0$. Then $\cW$ is flat over a
neighbourhood of $t=0$, $\cW_0=Y_1\cup\Theta$ as schemes, and the fibres $W_t$ for small $t$ are of pure dimension
$k$, contained in $\{F-t\det M=0\}$, of class the transport of $[Y_1]+(-1)^{c-1}[Y_2]+m'\eta^k$.
\end{proposition}
\begin{proof}[Sketch]
(i) Cofactor expansions; the last is Cramer's rule applied to the $n\times c$ matrix of the first $n$ columns of
$M$. (ii) As in Steps~2--3 of \S\ref{sec:R}, with the chain $V(p,h_n,h_0,..,h_{i-1},q_i,..,q_{n-1})$. (iii)
Outside $\Sigma$, $V(I)$ is locally cut out by $k+1$ equations: $(p,w)$ where $h_n\ne0$, and $(p,q_{<n},w_i)$ where
$f_i\neq0$ (by (i)); at $t=0$ these are the equations of $Y_1$, resp. $\Theta$, so $t$ is a non-zero-divisor there.
$\cW_0\subset V(I(0))$ agrees with it outside $\Sigma$, and $V(I(0))$ has no associated point in $\Sigma$. Purity follows
from a depth argument, and the class from (A).
\end{proof}

Once the sketch is completed this gives a proof of Corollary~4(d) for $\tau=0$ that does not use
Lemma~\ref{lem:NW}. For $c=2$, $W_t$ is a complete intersection in $V(p)$ (Hilbert--Burch; the determinantal
family of \cite{P2}); for $c=3$, a Pfaffian subscheme of codimension $3$ in $V(p)$ (Buchsbaum--Eisenbud; for
linear data, the Pfaffian cubics of \cite{P1} and, apparently, of Kloosterman \cite{Kl}). $W_t$ need not be the
closure of $\{p=0,\ f_jh_n=t\,a_j\}\smallsetminus\{h_n=0\}$: at special $X$ it has fixed components inside
$\{h_n=0\}$ (referee~11: skew lines on a quartic surface for which $M_{00},M_{10}$ have a common linear factor
modulo $h_1$). For $\tau>0$ the analogous statement with $a_j$ replaced by $Na_j$ is asserted in the working
notes and not proved here.

\section{A formula for \texorpdfstring{$A$}{A}}\label{sec:conj}

Assume $\tau=0$ and the range, and put $r:=d-\sigma=2k+2-(k-1)d>0$ (if $\sigma\le0$ then $K=0$). By Step~1 of
\S\ref{sec:G} and Lemma~\ref{lem:mu}, $K=I_1^\perp\cap I_2^\perp=(\Phi^1\cap\Phi^2)_\sigma=[(f)(h)+(g)+(q)+(p)+(s)]_\sigma$:
the degree-$\sigma$ part of the ideal whose degree-$d$ part is $T\VZ$. So every $Q\in K$ is the value of $d\mu$ on
a \emph{virtual deformation of weight $-r$} of the data, a tuple
$\dot v=(\dot f,\dot h,\dot p,\dot s,\dot M)$ of forms of the degrees of $f,h,p,s,M$ lowered by $r$ (zero where
this is negative):
\[
Q(\dot v)=f^\T\dot M\,h+\sum_i\dot f_ig_i+\sum_j\dot h_jq_j+\sum_l\dot p_ls_l+\sum_l\dot s_lp_l ,\qquad
Q(\dot v)\cdot N=d\mu(\dot vN)\in T\VZ\quad(N\in S_r).
\]
For $P\in S_d$ let
\begin{align*}
\mathfrak D_{\dot v}(P):={}&\sum_i\Res_S\Bigl[\frac{P\dot f_i}{(\dots,f_i^2,\dots)}\Bigr]+\sum_j\Res_S\Bigl[\frac{P\dot h_j}{(\dots,h_j^2,\dots)}\Bigr]\\
&+\sum_l\Res_S\Bigl[\frac{P\dot p_l}{(\dots,p_l^2,\dots)}\Bigr]+\sum_l\Res_S\Bigl[\frac{P\dot s_l}{(\dots,s_l^2,\dots)}\Bigr],
\end{align*}
where $(\dots,\phi^2,\dots)$ is the sequence $(f,p;h,s)$ with $\phi$ replaced by $\phi^2$ (all symbols are evaluated
in the socle degree), and
\[
D_{\dot v}:=\varrho\bigl(\operatorname{tr}(\dot M\cdot\operatorname{adj}M)\bigr)-\mathfrak D_{\dot v}(\det M).
\]
$D_{\dot v}$ is the formal derivative along $\dot v$ of the residue symbol $\Res_S[\det M/(f,p;h,s)]$ as a function
of the data. (The symbol itself is zero for degree reasons; its derivative along $\dot v$ has the right degree.)

\begin{theorem}[formula for $A$; read by referee 15]\label{thm:FA}
Assume $\tau=0$ and the range, and let $\kappa$ be the constant of Lemma~\ref{lem:PF}\textup{(4)}, so
that $\pi_1(Q'\det M)=\kappa\,\varrho(Q')$ for $Q'\in S_\sigma$; and
normalise the excess direction in \S\ref{sec:G} by $\omega:=u_1(\det M)$, i.e. $G_{\lz}=\det M$. Then for every
virtual deformation $\dot v$ of weight $-r$:
\begin{enumerate}
\item for all $P\in S_d$:\quad
$(k+1)\,\varphi_1\bigl(Q(\dot v)P\det M\bigr)=\pi_1\bigl(P\operatorname{tr}(\dot M\operatorname{adj}M)\bigr)+\kappa\,\mathfrak D_{\dot v}(P)$,\\
\phantom{for all $P\in S_d$:\quad}
$(k+1)\,\varphi_2\bigl(Q(\dot v)P\det M\bigr)=\pi_2\bigl(P\operatorname{tr}(\dot M\operatorname{adj}M)\bigr)-\lz\kappa\,\mathfrak D_{\dot v}(P)$;
\item $A(Q(\dot v))=\dfrac{k\kappa}{2}\,D_{\dot v}$.
\end{enumerate}
In particular $D_{\dot v}$ depends only on $Q(\dot v)$, and $A=0$ at $X$ if and only if $D_{\dot v}=0$ for all $\dot v$.
\end{theorem}
\begin{proof}
(1) Fix $\dot v$, $N\in S_r$, $Q'\in S_\sigma$, and move the data along the line
$(f,h,p,M,s)+\epsilon\,\dot vN$. By Lemma~\ref{lem:mu} the hypersurfaces $X_\epsilon$ lie in $\VZ$,
$F_\epsilon=F+\epsilon\,Q(\dot v)N+O(\epsilon^2)$, $\det M_\epsilon=\det M+\epsilon N\operatorname{tr}(\dot M\operatorname{adj}M)+O(\epsilon^2)$,
and $[Y_1(\epsilon)]$ is the transport of $[Y_1]$. By Lemma~\ref{lem:det} and Lemma~\ref{lem:PF}(1), (4) at $X_\epsilon$, with the
\emph{same} constant $\kappa$,
\begin{equation}\label{eq:first}
\pi_1^{X_\epsilon}(Q'\det M_\epsilon)=\kappa\,\Res_S\bigl[Q'/(f_\epsilon,p_\epsilon;h_\epsilon,s_\epsilon)\bigr].
\end{equation}
Differentiate at $\epsilon=0$. On the left, by (G), the derivative of a period of a fixed form over a flat class
in the direction $\dot F$ is $-(k+1)\varphi_1(\,\cdot\,\dot F)$, so we get
$-(k+1)\varphi_1(Q'\det M\cdot Q(\dot v)N)+\pi_1(Q'N\operatorname{tr}(\dot M\operatorname{adj}M))$. On the right we use
\[
\frac{d}{d\epsilon}\Big|_0\Res_S\Bigl[\frac{P'}{(\dots,\phi+\epsilon\psi,\dots)}\Bigr]=-\Res_S\Bigl[\frac{P'\psi}{(\dots,\phi^2,\dots)}\Bigr],
\]
which follows from the transformation law: $(\phi+\epsilon\psi)(\phi-\epsilon\psi)=\phi^2-\epsilon^2\psi^2$ gives
$\Res_S[P'/(\dots,\phi+\epsilon\psi,\dots)]=\Res_S[P'(\phi-\epsilon\psi)/(\dots,\phi^2-\epsilon^2\psi^2,\dots)]$, an
even function of $\epsilon$ minus $\epsilon$ times an even function with value $\Res_S[P'\psi/(\dots,\phi^2,\dots)]$ at
$0$. So the right side has derivative $-\kappa\,\mathfrak D_{\dot v}(Q'N)$. This is the first identity for $P=Q'N$;
both sides are linear in $P$ and the products $Q'N$ span $S_\sigma S_r=S_d$. For the second identity repeat with
$\pi_2$: by Lemma~\ref{lem:det} and Lemma~\ref{lem:PF}, $\pi_2(Q'\det M)=(-1)^c\kappa\,\varrho(Q')=-\lz\kappa\,\varrho(Q')$ (the
exchange of the blocks $f$ and $h$ is $c$ transpositions; this also follows from Theorem~1, which puts $\det M$
in $\ker(u_1+\lz u_2)$).

(2) Put $\omega:=u_1(\det M)=\kappa\varrho$; then $u_2(\det M)=-\lz\omega$. Choose $G_2\in T_1$ with $u_2(G_2)=\omega$
and put $G_1:=\det M+\lz G_2$. Then $u_1(G_1)=\omega$ and $u_2(G_1)=0$: an admissible choice in Step~4 of
\S\ref{sec:G}, with $G_{\lz}=G_1-G_2/\lz=\det M$. Let $Q=Q(\dot v)\in K$. By Theorem~1, $\Phi^{\delta_{\lz}}_Q$ vanishes
on the arc $F-t\det M$, so $\varphi_{\delta_{\lz}}(Q\det M^2)=0$, i.e. $A_0+B_0/\lz=0$; hence
$A_0(Q)=-\lz\varphi_2(QG_2^2)+2\lz\varphi_1(QG_1G_2)$. By the key lemma (Step~3: $\varphi_1(QG_2^2)=0$,
$\varphi_2(QG_1^2)=0$),
$\varphi_1(QG_1G_2)=\varphi_1(Q\det M\,G_2)$, and, writing $G_2=\lz(G_1-\det M)$,
$\varphi_2(QG_2^2)=-\varphi_2(Q\det M^2)-2\lz\varphi_2(Q\det M\,G_2)$. Therefore
\[
A_0(Q)=\lz\bigl[\varphi_2(Q\det M^2)+2\,\varphi_{\delta_{\lz}}(Q\det M\,G_2)\bigr],\qquad
\varphi_{\delta_{\lz}}=\varphi_1+\lz\varphi_2 .
\]
Now apply (1), with $\theta:=\operatorname{tr}(\dot M\operatorname{adj}M)\in S_\sigma$. For $P=G_2$ the terms with
$\mathfrak D$ cancel and $(k+1)\varphi_{\delta_{\lz}}(Q\det M\,G_2)=u_1(G_2)(\theta)+\lz u_2(G_2)(\theta)=\lz\kappa\varrho(\theta)$.
For $P=\det M$: $(k+1)\varphi_2(Q\det M^2)=-\lz\kappa[\varrho(\theta)+\mathfrak D_{\dot v}(\det M)]$. So
$(k+1)A_0(Q)=\kappa[\varrho(\theta)-\mathfrak D_{\dot v}(\det M)]=\kappa D_{\dot v}$, and $A=c_0A_0$ with
$c_0=k(k+1)/2$. Only Steps~1, 3 and~4 of \S\ref{sec:G} were used.
\end{proof}

The proof uses no pole-order reduction: every second-order quantity that occurs is of the form
$\varphi_i(Q\,P\det M)$ with $Q\in K$, and is obtained by differentiating the first-order identity
\eqref{eq:first} along $\VZ$ in the direction $Q N\in T\VZ$. It uses that $\kappa$ does not depend on $X$
(Lemma~\ref{lem:PF}(4)); alternatively, \eqref{eq:A2} follows from \eqref{eq:pw} and the same differentiation of the
algebraic identity of Lemma~\ref{lem:det}, at the single point $X$.

\begin{corollary}\label{cor:Anz}
Assume $\tau=0$ and the range, and consider a degree type $(k,d,c;\deg f;\deg h;\deg p)$ for which a
smooth $X$ exists and $K\ne0$. Then $A\ne0$ for general data $(f,h,p,M,s)$ of that type. Hence for general data,
and therefore at the general point $X'$ of $\VZ$ near any $X$ of the setting, $\NL(\delta_\lambda)$ is $\VZ$ with a
double structure for every $\lambda\notin\{0,\lz\}$.
\end{corollary}
\begin{proof}
The last sentence follows from the first because $\mu$ is a submersion onto $\VZ$ (Lemma~\ref{lem:mu}) and the
set of good data is Zariski-open and dense. Suppose first that $d_i+e_j\le\sigma$ for some $(i,j)$. It suffices
that $\operatorname{adj}(M)_{ji}\notin(f,h,p,s)$: then, the pairing on $\cA$ being perfect,
$\varrho(x^\alpha\operatorname{adj}(M)_{ji})\ne0$ for some monomial $x^\alpha$ of degree $\sigma-d_i-e_j$, and for
$\dot M=x^\alpha E_{ij}$ one has $Q(\dot v)=f_ih_jx^\alpha$ and $D_{\dot v}=\varrho(x^\alpha\operatorname{adj}(M)_{ji})$.
This is an open condition on the irreducible space of data, restricted to the open set where $(f,h,p,s)$ is a
regular sequence; so one instance in that set suffices (with $X$ possibly singular). Take $f,h,p,s$ general,
$\ell$ a general linear form, the complementary $(c-1)\times(c-1)$ submatrix of $M$ a monomial matrix whose
nonzero entries are powers of $\ell$ along a matching, and row $i$ and column $j$ of $M$ arbitrary. The cofactor
is then $\pm\ell^{d_i+e_j}$ (its degree is $d_i+e_j$ because $\tau=0$), which is not in $(f,h,p,s)$ for general
$\ell$, because powers of linear forms span $S_{d_i+e_j}$ and $\cA_{d_i+e_j}\ne0$. This needs a matching of the
submatrix along which all entry degrees are $\ge0$; that is automatic for $c=2$ and was checked by enumeration
(referees~15 and~16) for all degree types in the range with $k=2,3,4$, which are all types with $k\ge2$. For
$k=1$ and for $k=2,4$, $K\ne0$ implies that such $(i,j)$ exists ($k=1$: if $((f)(h))_\sigma=0$ then
$(g)_\sigma=(q)_\sigma=0$ as well; $k=2,4$: enumeration). For $k=3$ ($d=3$, $\sigma=1$) $K$ is spanned by one linear
form: $p$, or $s$ (if $p$ is a quadric), or, for $c=4$, the partner $q_j$ or $g_i$ of the quadric among $f,h$. The
seven linear forms among $f,h,p,s$ cut out a point $o$, and in coordinates in which they are coordinate
functions $D_{\dot v}=\mp\det M(o)/\phi(o)^2$, where $\phi$ is the quadric among $f,h,p,s$ ($\phi=s$, $p$, or the
quadric among $f,h$). Finally $\det M(o)\ne0$ for general $M$.
\end{proof}

\run{Checks and special cases.} (i) \emph{Numerical} (modulo primes; the second agent's tests of part (1) are
its own identities, built on the first agent's program module; files
\texttt{formula\_g/REPORT\_FG.md}, \texttt{referee15\_report.md}). Before the proof existed, the formula was found
by comparison with the second-order term computed by the program of \S\ref{sec:comp}: that quantity equals
$(-1)^{k+1}D_{\dot v}$ exactly, on every generator $f_ih_jx^\alpha$, $g_ix^\alpha$, $q_jx^\alpha$, $s_lx^\alpha$,
$p_lx^\alpha$ of $K$, in $25$ runs covering $12$ types ($k=1$: $(1,2)+(1,2)$ ($d=6$), $(1,1)+(3,3)$ ($d=8$), $(1,4)+(2,2)$ ($d=9$, not in Table~\ref{tab:types}); $k=2$,
$d=4$: all six types of Table~\ref{tab:types}; $k=2$, $d=5$: $(1,1)+(1,2)$ with $p$ a quadric, one instance;
$k=3$, $d=3$: the two types of Table~\ref{tab:types}, where $K$ is one-dimensional; the type with $p$ a quadric
and the mirror of the $c=4$ type were not run); in all, $247$ generator test vectors and no failure. According
to referee~15's translation of the program's quantity into $A$ (which is in its report, not in this manuscript),
the constant $(-1)^{k+1}$ is the ratio of $\kappa$ to the program's normalisation, as Theorem~\ref{thm:FA}
predicts. This formula was the first one tried; no alternative was tested. The right-hand side vanishes on all linear relations among the generators in the three types
where there are any. Part (1) of the theorem was tested separately, in translation to the Jacobian ring, in $7$
runs covering $6$ types ($k=1,2,3$, both signs of $\lz$; three random $P$ per test, one prime; not for $d=5$ or
$c\ge4$): no failure. (ii) For $k=1$, $c=2$ and $\omega=\varrho$
the theorem gives $A(f^\T\dot Mh)=\frac1{2\kappa}\varrho(\operatorname{tr}(\dot M\operatorname{adj}M))$: the formula
claimed in \cite{P2} under a hypothesis on one of the curves, here without that hypothesis and on
all of $K$.
(iii) For two $3$-planes meeting in a point $o$ on a cubic sixfold, $K=\langle p\rangle$ and
$D_{\dot v}=-\Res_S[\det M/(f,p;h,s^2)]$, which in coordinates where $f,h,p$ are coordinate functions is
$\pm\det M(o)/s(o)^2$: the factor $\det C$, $C=M(o)$, of Theorem~A of \cite{P1}. (iv) \emph{Where $A$ vanishes.}
$A=0$ iff $\operatorname{adj}(M)_{ji}\in(f,h,p,s)$ whenever $d_i+e_j\le\sigma$ and
$\Res_S[x^\alpha\det M/(\dots,\phi^2,\dots)]=0$ for every $\phi$ among the $f_i,h_j,p_l,s_l$ of degree $\ge r$ and every
monomial $x^\alpha$ of degree $\deg\phi-r$. Both hold when all entries of $M$ lie in $(f,h,p,s)$. Numerically: at
eight smooth instances satisfying both conditions (six types; in five of the eight instances, of four types,
all entries of $M$ were chosen in $(f,h,p,s)$) the computed second-order term is zero on all of $K$; at a smooth instance of type
$(1,1,1)+(1,2,2)$, $d=4$, satisfying only the first, it vanishes on $((f)(h))_\sigma$ and not on $(q)_\sigma$.

\section{Several excess directions}\label{sec:FA2}

Assume $\tau\ge0$ and the range. Put $\sigma':=kd-2k-2=\sigma-\tau$ (the degree of the forms $Q$ of
\S\ref{sec:G}), $r:=d-\sigma'$ (the range says $r>\tau$), $E:=\cA_\tau$, and for $N\in S_\tau$ let
$\omega_N:=u_1(N\det M)\in I_1\cap I_2$; by Lemma~\ref{lem:det}, $N\mapsto\omega_N$ induces an isomorphism
$E\cong I_1\cap I_2$. As in \S\ref{sec:conj}, $K=(I_1+I_2)^\perp=[(f)(h)+(g)+(q)+(p)+(s)]_{\sigma'}$ consists of the
values $Q(\dot v)$ of $d\mu$ on the virtual deformations $\dot v$ of weight $-r$, and $\mathfrak D_{\dot v}(P)$ is
defined as there, now for $P\in S_{d+\tau}$. Let $c_0:=k(k+1)/2$.

\begin{theorem}[read by referee 17]\label{thm:FA2}
Assume $\tau\ge0$ and the range.
\begin{enumerate}
\item \emph{Structure.} For every $\lambda\ne0$, $\NL(\delta_\lambda)$ lies in a smooth germ $\cM_\lambda$ of
dimension $\dim\VZ+\et$, in which $\VZ=\{y_1=\dots=y_{\et}=0\}$; for $Q\in K$ the function
$\phi_Q:=\Phi^{\delta_\lambda}_Q|_{\cM_\lambda}$ lies in the square $(y)^2$ of the ideal of $\VZ$; and
$\NL(\delta_\lambda)=\{\phi_Q=0,\ Q\in K\}\subset\cM_\lambda$.
\item \emph{Law.} Let $G_\lambda\colon E\to T\NL(\delta_\lambda)/T\VZ$ be the isomorphism with
$u_1(G_\lambda(N))=\omega_N$ (so $G_{\lz}(N)=N\det M$). There is a symmetric bilinear map
$\mathcal Q_0\colon E\times E\to K^*$, independent of $\lambda$, such that for $Q\in K$, $N,N'\in E$ and
$\lambda\ne0$
\[
\tfrac12\Hess\phi_Q\bigl(G_\lambda(N),G_\lambda(N')\bigr)=c_0\,\frac{\lambda-\lz}{\lambda}\,\mathcal Q_0(N,N')(Q).
\]
\item \emph{Formula.} For every virtual deformation $\dot v$
of weight $-r$ and $N,N'\in S_\tau$:
\[
\mathcal Q_0(N,N')\bigl(Q(\dot v)\bigr)=\frac{\kappa}{k+1}\Bigl[\varrho\bigl(NN'\operatorname{tr}(\dot M\operatorname{adj}M)\bigr)
-\mathfrak D_{\dot v}\bigl(NN'\det M\bigr)\Bigr],
\]
the formal derivative along $\dot v$ of $\Res_S[NN'\det M/(f,p;h,s)]$ with $N,N'$ held fixed.
\item \emph{Consequence.} If the quadrics $N\mapsto\mathcal Q_0(N,N)(Q)$, $Q\in K$, have no common zero in
$\PP(E)$ at $X$, then for every $\lambda\notin\{0,\lz\}$: $\NL(\delta_\lambda)_{\mathrm{red}}=\VZ$ as germs at $X$, and
$\NL(\delta_\lambda)$ is non-reduced at every point of $\VZ$ near $X$.
\end{enumerate}
For $\et=1$, parts (1), (2), (4) are Corollary~4 with $A=c_0\mathcal Q_0(N,N)$, $N$ a generator of $E$; for
$\et=1$ and $\tau=0$, part (3) is Theorem~\ref{thm:FA}.
\end{theorem}
\begin{proof}
(1) Choose $Q_1,\dots,Q_m$ spanning a complement of $K$ in $S_{\sigma'}$. The $\Phi^{\delta_\lambda}_{Q_a}$ have
independent differentials at $X$, and $m=\codim T\NL(\delta_\lambda)$; let $\cM_\lambda$ be their common zero germ,
smooth of dimension $\dim\VZ+\et$ (Theorem~2(2)). As $\Phi_Q$ is linear in $Q$,
$\NL(\delta_\lambda)=\{\phi_Q=0,\ Q\in K\}$ in $\cM_\lambda$, and $\phi_Q$ vanishes on $\VZ$. Let $X'\in\VZ$ be near
$X$. The Zariski tangent space of $\NL(\delta_\lambda)$ at $X'$ is $\{v\in T_{X'}\cM_\lambda:\ d\phi_Q(X')(v)=0,\ Q\in K\}$;
by (G), (A) and Theorem~2(2) at $X'$ (the setting holds there by Lemma~\ref{lem:mu}, and the range and $\et$
depend only on the degrees) it has dimension $\dim\VZ+\et=\dim T_{X'}\cM_\lambda$. So $d\phi_Q(X')=0$ for all
$Q\in K$ (the space $K$ is that of $X$; that $K$ varies with $X'$ plays no role). A function on a smooth germ
vanishing with its differential along the smooth subgerm $\{y=0\}$ lies in $(y)^2$.

(2) Choose a linear map $N\mapsto G_2(N)\in T_1$ with $u_2(G_2(N))=\omega_N$ (possible since $u_2(T_1)=I_2$ in
the range) on a complement of $(f,h,p,s)_\tau$ in $S_\tau$, extended by $0$ on $(f,h,p,s)_\tau$, and put $G_1(N):=N\det M+\lz G_2(N)$,
$G_\lambda(N):=G_1(N)-G_2(N)/\lambda$. By Lemma~\ref{lem:det}, $u_1(G_1(N))=\omega_N$ and $u_2(G_1(N))=0$, so
$G_1(N)\in T_2$, $(u_1+\lambda u_2)(G_\lambda(N))=0$, $u_1(G_\lambda(N))=\omega_N$, and $G_{\lz}(N)=N\det M$. By (1)
the form $\tfrac12\Hess\phi_Q$ on $T\NL(\delta_\lambda)$ has $T\VZ$ in its kernel, and the class of $G_\lambda(N)$
modulo $T\VZ$ is determined by $u_1(G_\lambda(N))=\omega_N$; so the left side of the law does not depend on the
choices. By (G) it equals $c_0\varphi_{\delta_\lambda}(QG_\lambda(N)G_\lambda(N'))$ (the differential of
$\Phi_Q$ vanishes at $X$ for $Q\in K$). Expand $\varphi_{\delta_\lambda}=\varphi_1+\lambda\varphi_2$ and use the key
lemma of \S\ref{sec:G} (its Step~3 uses only hypotheses (a) and (b), which hold in the range) in polarised form ($\varphi_1(QGG')=0$ for $G,G'\in T_1$, $\varphi_2(QGG')=0$ for
$G,G'\in T_2$): with $G_i=G_i(N)$, $G_i'=G_i(N')$,
\begin{align*}
\varphi_{\delta_\lambda}\bigl(QG_\lambda(N)G_\lambda(N')\bigr)&=A_0+B_0/\lambda,\\
A_0&:=\varphi_1(QG_1G_1')-\varphi_2(QG_1G_2')-\varphi_2(QG_2G_1'),\\
B_0&:=\varphi_2(QG_2G_2')-\varphi_1(QG_1G_2')-\varphi_1(QG_2G_1').
\end{align*}
By Theorem~1, $\Phi^{\delta_{\lz}}_Q$ vanishes on the affine space $F+\det M\cdot S_\tau$ near $F$; its second-order
part there is $c_0\varphi_{\delta_{\lz}}(QN\det M\,N'\det M)$, so $A_0+B_0/\lz=0$ and
$\mathcal Q_0:=A_0$ satisfies the law.

(3) \emph{Step A.} Fix $\dot v$, $N\in S_\tau$, $N_r\in S_r$, $Q'\in S_{\sigma'}$ and move the data along
$(f,h,p,M,s)+\epsilon\,\dot vN_r$, as in the proof of Theorem~\ref{thm:FA}. By Lemma~\ref{lem:det} and Lemma~\ref{lem:PF},
$\pi_1^{X_\epsilon}(Q'N\det M_\epsilon)=\kappa\Res_S[Q'N/(f_\epsilon,p_\epsilon;h_\epsilon,s_\epsilon)]$; differentiating
at $\epsilon=0$ exactly as there, and using $S_{\sigma'}S_r=S_d$, we get for all $P\in S_d$, $N\in S_\tau$, with
$\theta:=\operatorname{tr}(\dot M\operatorname{adj}M)$:
\begin{equation}\label{eq:A2}
\begin{aligned}
(k+1)\varphi_1\bigl(Q(\dot v)PN\det M\bigr)&=\pi_1(PN\theta)+\kappa\,\mathfrak D_{\dot v}(PN),\\
(k+1)\varphi_2\bigl(Q(\dot v)PN\det M\bigr)&=\pi_2(PN\theta)-\lz\kappa\,\mathfrak D_{\dot v}(PN).
\end{aligned}
\end{equation}
\emph{Step B.} Let $Q=Q(\dot v)$. From $A_0=-\lz B_0$, the key lemma (which gives
$\varphi_1(QG_1G_2')=\varphi_1(QN\det M\,G_2')$ and $\varphi_2(QG_1G_1')=0$) and $G_2=\lz(G_1-N\det M)$:
\[
A_0=\lz\Bigl[\varphi_2\bigl(QNN'\det M^2\bigr)+\varphi_{\delta_{\lz}}\bigl(QN\det M\,G_2'\bigr)
+\varphi_{\delta_{\lz}}\bigl(QN'\det M\,G_2\bigr)\Bigr].
\]
\emph{Step C.} By \eqref{eq:A2} with $P=G_2'$ (the $\mathfrak D$-terms cancel, $\kappa-\lz\cdot\lz\kappa=0$):
$(k+1)\varphi_{\delta_{\lz}}(QN\det M\,G_2')=u_1(G_2')(N\theta)+\lz u_2(G_2')(N\theta)=\lz\omega_{N'}(N\theta)
=\lz\kappa\varrho(NN'\theta)$, and the same for the third term. By \eqref{eq:A2} with $P=N'\det M$:
$(k+1)\varphi_2(QNN'\det M^2)=\pi_2(N'\det M\,N\theta)-\lz\kappa\mathfrak D_{\dot v}(NN'\det M)
=-\lz\kappa[\varrho(NN'\theta)+\mathfrak D_{\dot v}(NN'\det M)]$. Hence
$(k+1)A_0=\kappa[\varrho(NN'\theta)-\mathfrak D_{\dot v}(NN'\det M)]$. Both sides of the formula vanish for
$N\in(f,h,p,s)_\tau$: the right side because a generator $\phi$ of degree $\le\tau$ has $\dot\phi=0$ ($r>\tau$).

(4) In coordinates $(x',y)$ on $\cM_\lambda$, $\phi_Q=\sum_{a\le b}y_ay_b\upsilon_{Q,ab}(x',y)$ by (1), and by (2)
the quadratic form $\sum\upsilon_{Q,ab}(X,0)y_ay_b$ is $c_0\frac{\lambda-\lz}{\lambda}\mathcal Q_0(N,N)(Q)$ in the
coordinates of $E$ given by $dy\circ G_\lambda$. Let $Q_1,\dots,Q_t$ be a basis of $K$. If the quadrics have no
common zero, then by compactness of the unit sphere and continuity there are $c>0$ and a neighbourhood of $X$ in
$\VZ$ on which $\max_a|\sum\upsilon_{Q_a,\cdot}(x',0)yy|\ge c|y|^2$, while $\phi_{Q_a}$ differs from its quadratic
part by $O(|y|^3)$ uniformly. So the common zero set of the $\phi_Q$ is $\{y=0\}$ near $X$. The Zariski tangent
space of $\NL(\delta_\lambda)$ at points of $\VZ$ has dimension $\dim\VZ+\et>\dim\VZ$.
\end{proof}

\begin{corollary}[one excess direction on surfaces, $\tau>0$; read by referee 18]\label{cor:surf}
Consider one of the four surface types listed after Corollary~4 ($k=1$, $\tau>0$,
$\et=1$, the range, $K\ne0$). Exactly one of the forms $f,h,p,s$ is not linear; call it $\phi$, let $\phi^*$ be
its partner in \eqref{eq:F} (so $F\equiv\phi\phi^*$ modulo the other products), and let $o$ be the point where the
three linear forms vanish. Work in coordinates of $\PP^3$ in which the three linear forms and a fourth linear
form $y$ are the coordinate functions (so $y(o)=1$; in other coordinates the right side below is multiplied by
$J^{-1}$, $J$ the determinant of these four linear forms), and let $N=y^\tau$ generate $E$. Then, with $A$
normalised by $G_{\lz}=N\det M$ and $r=4$, for every virtual deformation $\dot v$ of weight $-4$:
\[
A\bigl(Q(\dot v)\bigr)=\pm\frac{\kappa}{2}\Bigl[\frac{\operatorname{tr}(\dot M\operatorname{adj}M)(o)}{\phi(o)}
-\frac{\det M(o)\,\dot\phi(o)}{\phi(o)^2}\Bigr],
\]
the sign depending only on the type. Consequently:
\begin{enumerate}
\item if $d\ge6$, then $A\ne0$ at every smooth $X$ of the type;
\item if $d=5$ (the first three types), then $K=\langle\phi^*\rangle$ and $A\ne0$ if and only if $\det M(o)\ne0$,
which holds for general data. For two coplanar lines on a quintic this says: the residual plane cubic
$\{M_{00}=0\}$ does not pass through the point where the lines meet.
\end{enumerate}
\end{corollary}
\begin{proof}
In these coordinates $\cA=\CC[y]/(y^a)$ with $a=\deg\phi$, and by the transformation law every symbol
$\Res_S[P/(\dots)]$ occurring is, up to one common sign (the position of $\phi$ in the sequence), the
one-variable residue of the restriction to the line $x_0=x_1=x_2=0$: $\varrho(P)=\pm P_{a-1}/\phi(o)$
and $\Res_S[P/(\dots,\phi^2,\dots)]=\pm P_{2a-1}/\phi(o)^2$, where $P_j$ is the coefficient of $y^j$ in the
restriction ($\phi(o)\ne0$ by Lemma~\ref{lem:alg}). A linear form $\ell$ among $f,h,p,s$ has $\dot\ell=0$, since
$r\ge2$. Here $\sigma'=d-4$, $r=4$, and $\deg\phi\ge4$ in all four types ($d-1$ with $d\ge5$, or $d-2$ with
$d\ge6$). Insert in Theorem~\ref{thm:FA2}(3) with $N=N'=y^\tau$, $k=1$, $c_0=1$; the degrees are such that
$y^{2\tau}\operatorname{tr}(\dot M\operatorname{adj}M)$ has degree $a-1$ and $y^{2\tau}\det M\,\dot\phi$ degree $2a-1$.

(1) (Referee~18.) For $d\ge6$ one entry $\dot M_{i0}$ of $\dot M$ may be taken to be $y^{d-6}$, with the other
entries of $\dot M$, and $\dot\phi$, zero. For $c=1$ (two coplanar lines; two plane curves of degree $d-1$)
$\operatorname{adj}M=(1)$ and $A(f_0h_0y^{d-6})=\pm\frac\kappa2/\phi(o)\ne0$. For $c=2$,
$\operatorname{tr}(\dot M\operatorname{adj}M)=\dot M_{00}M_{11}-\dot M_{10}M_{01}$, so it suffices that
$(M_{01}(o),M_{11}(o))\ne0$: for the line and the curve of degree $d-1$ these are constants, not both zero
because $\det M\not\equiv0$; for the curve of degree $d-2$, $\nabla F(o)=h_1(o)\bigl(M_{01}(o)\nabla f_0+M_{11}(o)\nabla f_1\bigr)$,
which is nonzero as $X$ is smooth at $o$. (2) For $d=5$ only $\dot\phi\in\CC$ is nonzero, $K=\langle\phi^*\rangle$,
and $A(\phi^*)=\mp\frac\kappa2\det M(o)/\phi(o)^2$; the condition $\det M(o)\ne0$ is open and non-empty on the
irreducible space of data, hence compatible with smoothness of $X$. For coplanar lines $\det M=M_{00}$, and the
plane section of $X$ is $\{f_0M_{00}h_0=0\}$.
\end{proof}

\begin{corollary}[cubic sixfolds, two $3$-planes meeting in a line]\label{cor:six}
Let $X\subset\PP^7$ be a smooth cubic containing two $3$-planes $\Pi_1=V(f_0,f_1,p_1,p_2)$,
$\Pi_2=V(h_0,h_1,p_1,p_2)$ with $L=\Pi_1\cap\Pi_2$ a line ($k=d=3$, $c=2$, $\tau=1$, $\lz=-1$, $\et=2$,
$K=\langle p_1,p_2\rangle$), and let $\bar m$, $\bar s_1$, $\bar s_2$ be the restrictions of $\det M$, $s_1$, $s_2$ to
$L$ (binary quadrics).
\begin{enumerate}
\item The two quadrics of Theorem~\ref{thm:FA2}(4) are nonzero multiples of
$N\mapsto\Res_{\CC[y_1,y_2]}[\bar N^2\bar m/(\bar s_l^2,\bar s_{l'})]$, $l=1,2$, and they have no common zero if
and only if $\bar m$ has two distinct zeros on $L$ and $\bar m\notin\langle\bar s_1,\bar s_2\rangle$. In that case,
for every $\lambda\notin\{0,-1\}$, $\NL([\Pi_1]+\lambda[\Pi_2])_{\mathrm{red}}=\VZ$ near $X$, with a non-reduced
structure.
\item At every such $X$ and for every $\lambda\notin\{0,-1\}$ the germ of $\NL([\Pi_1]+\lambda[\Pi_2])$ at $X$ is
non-reduced, hence not smooth; for $\lambda=-1$ it is smooth (Theorem~2(4)).
\end{enumerate}
\end{corollary}
\begin{proof}
(1) $g_i,q_j,s_l,f_ih_j$ have degree $2>\sigma'=1$, so $K=\langle p_1,p_2\rangle$ and the only virtual deformations
are $\dot s_l\in\CC$; Theorem~\ref{thm:FA2}(3) gives $\mathcal Q_0(N,N)(p_l)=-\frac\kappa4\Res_S[N^2\det M/(f,p;h,\dots,s_l^2,\dots)]$,
and the transformation law reduces the symbol to $L$. As $X$ is smooth along $L$, the pencil
$\langle\bar s_1,\bar s_2\rangle$ has no base point; it contains exactly two squares (a pencil tangent to the conic of squares at $\ell^2$ consists
of multiples of $\ell$ and has a base point), so we may assume
$\bar s=(y_1^2,y_2^2)$. For $\bar m=\alpha y_1^2+\beta y_1y_2+\gamma y_2^2$ and $\bar N=ay_1+by_2$ the two quadrics
are $a(a\beta+2b\alpha)$ and $b(2a\gamma+b\beta)$ up to constants, with a common zero iff $\beta=0$ or
$\beta^2=4\alpha\gamma$ (referee~17). (2) The condition in (1) holds for general data (for instance on the Fermat
cubic with the planes $x_0+x_1=x_2+x_3=x_4+x_5=x_6+x_7=0$ and $x_0+x_2=x_1+x_3=x_4+x_6=x_5+x_7=0$: referee~17), and
$\mu$ is a submersion onto $\VZ$ (Lemma~\ref{lem:mu}); so $\NL$ is non-reduced at a dense set of points of $\VZ$
near $X$, and the non-reduced locus of an analytic space is closed.
\end{proof}

\run{Remarks.} (i) At special $X$ the quadrics can vanish identically: on the Fermat cubic with two planes of the
\emph{same} matching ($f=(x_0+x_1,x_2+x_3)$, $h=(x_0+\omega x_1,x_2+\omega x_3)$) one has $M|_L=0$; there (2) holds
but nothing is said about the support. This is the vanishing at the Fermat point met in \cite{P1,P2}.
(ii) As digested, \cite[Conj.~6, 7]{Mo} assert for $n=6$ that the locus is smooth exactly for the coefficients
$(1,-1)$, at the Fermat point, and report a computer proof of the non-smooth half for $n=6,8$;
Corollary~\ref{cor:six}(2) with Theorem~2(4) gives that dichotomy at every smooth $X$ for $n=6$. The comparison
with the wording of the conjectures has not been made by a person. The case was treated by direct computation
in \cite{P1}. (iii) Theorem~\ref{thm:FA2}(4) needs $\dim K\ge\et$; when $\dim K<\et$ the quadrics always have
common zeros and the structure of $\NL(\delta_\lambda)$ is not determined here. (iv) On the Fermat surface of
degree $d$ the $d$ lines in a plane $x_0=\zeta x_1$ are concurrent, so for two of them $M_{00}(o)=0$: by
Corollary~\ref{cor:surf}, $A=0$ on the Fermat quintic and $A\ne0$ for $d\ge6$ (both confirmed modulo a prime by
referee~18).

\run{Checks} (modulo primes; file \texttt{formula\_g2/REPORT\_FG2.md}; the program module is the one used for
\S\ref{sec:conj}, so these tests are not independent of the earlier ones). The law (2) and the formula (3) were
tested at $41$ instances of $17$ degree types in the range with $\tau\in\{1,2\}$, $\et\in\{1,2,3,6\}$,
$k\le3$, $c\le4$ and both signs of $\lz$ (three of the instances were special, the others random): the computed bilinear form
vanishes at $\lz$, the rescaled form is the same for five values of $\lambda\ne\lz$, and (3) holds exactly, with
the same normalisation constant as in \S\ref{sec:conj}, on $1222$ test triples $(N,N',Q(\dot v))$ ($1219$ of them
nonzero), including pairs $N\ne N'$ and all five kinds of generators of $K$. This includes coplanar lines for
$d=5,6,7$, i.e.\ the input of Corollary~\ref{cor:surf}; its closed form was tested afterwards by referee~18, modulo one
prime, in $30$ runs covering the four types (proportionality on up to $30$ generators per run, not the absolute
constant), including instances with $\det M(o)=0$ and the Fermat surfaces of degree $5,6,7$ (files
\texttt{referee18/}).
For the type of Corollary~\ref{cor:six} the two quadrics had nonzero resultant in $5$ of $5$ generic instances
(three primes) and a common zero in three instances built to violate the criterion ($\bar m$ a square;
$\det M\in(f,h,p,s)$). Exact over $\QQ$: the two-variable computation of Corollary~\ref{cor:six}(1) at six random
instances and at the Fermat instance of its proof (files \texttt{fa2\_check/}, \texttt{referee17/}). The
algebra of the normal form in the proof of Corollary~\ref{cor:six} is certified by the Lean~4 kernel (file
\texttt{SixfoldCriterionCert.lean}, ten theorems, nine over an arbitrary commutative ring and one over $\ZZ$: the two coefficients of
$\bar N^2\bar m$; identities exhibiting $\beta(\beta^2-4\alpha\gamma)a^3$ and $\beta(\beta^2-4\alpha\gamma)b^3$ as
combinations of the two quadrics; the common zeros $(1,0)$ if $\beta=0$ and $(2\alpha,-\beta)$ if
$\beta^2=4\alpha\gamma$; and, for one four-variable block of the Fermat cubic, the identity $x_0^3+\dots+x_3^3=f_0M_{00}h_0+p_1s_1$ with
$M_{00}|_L=6x_0$, $s_1|_L=3x_0^2$, whence $\bar m=36y_1y_2$, $\bar s=(3y_1^2,3y_2^2)$); it certifies nothing about
residues or Hodge loci. Parts (1)
and (4) of Theorem~\ref{thm:FA2} were not tested beyond first-order dimensions.


\subsection{Outside the range}\label{sec:FA3}

Without the range, $S_d/(T_1+T_2)=\cA_d$ need not vanish. Put
\[
E':=(0:_\cA\det M)_\tau=\ker\bigl(\cdot\det M\colon\cA_\tau\to\cA_d\bigr)\subset E,\qquad
\text{(S):}\quad \cdot\det M\colon\cA_\tau\to\cA_d\ \text{is surjective}.
\]
By Theorem~2(5), $T\NL(\delta_\lambda)/T\VZ\cong E'$ for $\lambda\notin\{0,\lz\}$; in the range $E'=E$ and (S) holds
trivially. Let $K:=(I_1+I_2)^\perp=[(f)(h)+(g)+(q)+(p)+(s)]_{\sigma'}$ as before (Lemma~\ref{lem:mu}) and
$K_\lambda:=(\operatorname{im}(u_1+\lambda u_2))^\perp\supseteq K$, the space of forms $Q$ with
$d\Phi^{\delta_\lambda}_Q(X)=0$; $r=d-\sigma'$ may now have any sign.

\begin{theorem}[read by referee 19]\label{thm:FA3}
In the setting, with $\tau\ge0$ and no range hypothesis, let $\lambda\notin\{0,\lz\}$.
\begin{enumerate}
\item[\textup{(0)}] $\dim K_\lambda-\dim K=\dim\operatorname{coker}(\cdot\det M\colon\cA_\tau\to\cA_d)$; so
$K_\lambda=K$ if and only if \textup{(S)} holds.
\item \emph{Law.} For $N\in E'$ let $G_\lambda(N)\in\ker(u_1+\lambda u_2)$ be the class modulo $T\VZ$ with
$u_1(G_\lambda(N))=\omega_N$. There is a symmetric bilinear map $\mathcal Q_0\colon E'\times E'\to K^*$,
independent of $\lambda$, such that for $Q\in K$ and $N,N'\in E'$
\[
\tfrac12\Hess\Phi^{\delta_\lambda}_Q\bigl(G_\lambda(N),G_\lambda(N')\bigr)=c_0\,\frac{\lambda-\lz}{\lambda}\,\mathcal Q_0(N,N')(Q),
\]
and $T\VZ$ lies in the kernel of this Hessian on $\ker(u_1+\lambda u_2)$.
\item \emph{Formula.} Assume $r\ge0$. Then the formula of Theorem~\ref{thm:FA2}(3) holds for
$N,N'\in S_\tau$ with classes in $E'$.
\item \emph{Not smooth.} If $\mathcal Q_0\ne0$ at $X$, then $\NL(\delta_\lambda)$ is not smooth at $X$ (as an
analytic space). If $\mathcal Q_0\ne0$ at one smooth $X$ of a degree type at which $\dim E'$ is minimal (or if,
for $r\ge0$, the right side of (2) is nonzero at one point of the space of data at which
$(f,h,p,s)$ is regular and $\dim E'$ is minimal, $X$ possibly singular, and the type admits a smooth $X$), then
$\NL(\delta_\lambda)$ is singular at every smooth $X$ of that type.
\item \emph{Structure.} Assume \textup{(S)} at $X$. Then the statements of Theorem~\ref{thm:FA2}(1) and (4) hold
near $X$ with $E$ replaced by $E'$: $\NL(\delta_\lambda)=\{\phi_Q=0,\ Q\in K\}$ in a smooth germ of dimension
$\dim\VZ+\dim E'$, $\phi_Q\in(y)^2$, and if the quadrics $N\mapsto\mathcal Q_0(N,N)(Q)$, $Q\in K$, have no common
zero in $\PP(E')$, then $\NL(\delta_\lambda)_{\mathrm{red}}=\VZ$ and $\NL(\delta_\lambda)$ is non-reduced along $\VZ$.
\end{enumerate}
\end{theorem}
\begin{proof}
The inputs are range-free: Lemmas~\ref{lem:T}, \ref{lem:mu}, \ref{lem:det}, Theorem~1, and Theorem~2(1), (4),
(5) and the embedding in (2) (referee~19 checked this in the proofs above). (0) $I_i\cong S_d/T_i$,
$\dim(I_1\cap I_2)=\et$ and $\dim(T_1\cap T_2)=\dim T_1+\dim T_2-\dim S_d+\dim\cA_d$ give
$\dim(I_1+I_2)-\operatorname{rank}(u_1+\lambda u_2)=\dim\cA_d-(\et-\dim E')$, by Theorem~2(5).

(1) The key lemma holds for $Q\in K=I_1^\perp\cap I_2^\perp$: $V_i$ is smooth with tangent space $T_i$,
$\Phi^{\delta_i}_Q$ vanishes on it and has zero differential at $X$. For $N\in S_\tau$ with
$N\det M=H_1+H_2$, $H_i\in T_i$, put $G_2(N):=-\lz H_1\in T_1$ and $G_1(N):=H_2=N\det M+\lz G_2(N)\in T_2$; then
$u_2(G_2(N))=\omega_N=u_1(G_1(N))$ by Lemma~\ref{lem:det}, and $G_\lambda(N):=G_1(N)-G_2(N)/\lambda$ has the
required properties. For $z\in T_1\cap T_2$ the key lemma gives
$\varphi_{\delta_\lambda}(QzG_\lambda(N))=\varphi_1(QzG_1)-\varphi_2(QzG_2)=\varphi_{\delta_{\lz}}(Qz\,N\det M)$, which
vanishes because $\NL(\delta_{\lz})$ is smooth with tangent space $T\VZ+\det M\cdot S_\tau$ (Theorem~2(4)),
$\Phi^{\delta_{\lz}}_Q$ vanishes on it, and $Q\in K\subset K_{\lz}$. Also $\varphi_i(Qzz')=0$ for
$z,z'\in T_1\cap T_2$ by the key lemma, and every element of $\ker(u_1+\lambda u_2)$ is $G_\lambda(N)+z'$ with
$N\in E'$ (Theorem~2(5)). So $T\VZ$ is in the kernel, and the form
depends only on the classes; another decomposition, or another representative of $N$ modulo $(f,h,p,s)_\tau$,
changes $G_\lambda(N)$ by an element of $T_1\cap T_2$. The expansion into $A_0+B_0/\lambda$ and the identity
$A_0+B_0/\lz=0$ are as in the proof of Theorem~\ref{thm:FA2}(2).

(2) Steps A--C of the proof of Theorem~\ref{thm:FA2}(3) use the range only through the existence of $G_2(N)$,
which is the condition $N\in E'$, and through $S_{\sigma'}S_r=S_d$, which needs $r\ge0$. (Both sides vanish for
$N\in(f,h,p,s)_\tau$ and $N'\in E'$: on the right, the term of $\mathfrak D_{\dot v}$ with $\phi$ squared is
$\varrho(a\,\dot\phi\,N'\det M)=0$ for $N=a\phi$.)

(3) Let $\cM_\lambda$ be the zero germ of the $\Phi_{Q_a}$, $Q_a$ spanning a complement of $K_\lambda$: smooth,
with tangent space $T\NL(\delta_\lambda)$. If $\NL(\delta_\lambda)$ were smooth it would equal $\cM_\lambda$, so
every $\Phi_Q|_{\cM_\lambda}$ would vanish and so would its Hessian, which for $Q\in K$ is the left side of the
law. For the second statement (referee~19): where $\dim E'$ is minimal, $K$ and $E'$ are vector bundles on the
connected space of data and $\mathcal Q_0$ is a holomorphic section, so it is nonzero off a proper analytic
subset; $\mu$ is open onto $\VZ$ (Lemma~\ref{lem:mu}), and the singular locus of an analytic space is closed.
For $r\ge0$ the right side of (2) is a regular function on the irreducible bundle of triples (data, $N$, $N'$)
with $N,N'\in E'$ over the open set where $(f,h,p,s)$ is regular and $\dim E'$ is minimal; so ``nonzero'', (S)
and ``no common zero'' at one point of it, even one with $X$ singular, hold on a dense open set, which meets
the open set where $X$ is smooth.

(4) (S) is open, so near $X$ on $\VZ$ the dimension of $E'$ is constant and $K_\lambda=K$ by (0); now argue as
for Theorem~\ref{thm:FA2}(1), (4).
\end{proof}

\begin{propallr}[the formula for every $r$; added in version 20; first proof read by referee 24, who found a
gap and supplied step (ii) below; present proof read by referee 25]
Assume $\tau\ge0$; the integer $r=d-\sigma'$ may have any sign. For $r<0$ call a \emph{virtual deformation of weight
$-r$} a tuple $\dot v=(\dot f,\dot h,\dot p,\dot s,\dot M)$ of forms of degrees $\deg f_i-r$, $\deg h_j-r$,
$\deg p_l-r$, $\deg s_l-r$ and $d-\deg f_i-\deg h_j-r$ for $\dot M_{ij}$ (zero if this number is negative; $\dot M_{ij}$
may be nonzero also where $M_{ij}=0$ for degree reasons), and define $Q(\dot v)$,
$\theta=\operatorname{tr}(\dot M\operatorname{adj}M)$ and $\mathfrak D_{\dot v}$ by the formulas of \S\ref{sec:conj}. Then
$K=\{Q(\dot v)\}$; the two identities \eqref{eq:A2} hold for all $P\in S_d$, $N\in S_\tau$; and the formula of
Theorem~\ref{thm:FA2}\textup{(3)} holds for all $N,N'\in S_\tau$ with classes in $E'$. So the hypothesis $r\ge0$ in
Theorem~\ref{thm:FA3}\textup{(2)} and \textup{(3)} is not needed.
\end{propallr}
\begin{proof}
Let $r<0$. By Lemma~\ref{lem:mu}, $K=(I_1+I_2)^\perp$ is the degree-$\sigma'$ part of $(f)(h)+(g)+(q)+(p)+(s)$
(\S\ref{sec:FA3}), and the multipliers of $f_ih_j$, $g_i$, $q_j$, $s_l$, $p_l$ in degree $\sigma'=d-r$ have the
degrees prescribed for $\dot M_{ij}$, $\dot f_i$, $\dot h_j$, $\dot p_l$, $\dot s_l$; so $K=\{Q(\dot v)\}$. Both sides of
\eqref{eq:A2} are linear in $\dot v$, so it suffices to treat (i) tuples with $\dot M=0$ and (ii) tuples whose only
nonzero entry is one $\dot M_{ij}=m$.

(i) The members $f_i,h_j,p_l,s_l$ have positive degree, and a form of degree $b-r$, $b\ge0$, is a sum
$\sum_aL_aw_a$ with $L_a$ running over the monomials of degree $-r$ and $w_a\in S_b$. Doing this for every entry,
$\dot v=\sum_aL_a\,w^{(a)}$, where each $w^{(a)}$ is a tuple of forms of the degrees of the data themselves, with zero
matrix part. The first part of Step~A of the proof of Theorem~\ref{thm:FA2}(3), for the family
$(f,h,p,M,s)+\epsilon\,w^{(a)}$ (a family of data, so no factor $N_r$ is needed) and for $Q'\in S_{\sigma'}$,
$N\in S_\tau$, gives
\[
(k+1)\varphi_1\bigl(Q(w^{(a)})\,Q'N\det M\bigr)=\kappa\,\mathfrak D_{w^{(a)}}(Q'N),
\]
and the same with $\varphi_2$ and $-\lz\kappa$ ($\theta=0$ here). $\mathfrak D_w(P')$ is defined for such a tuple $w$
and $P'\in S_{\sigma'+\tau}$ by the formula of \S\ref{sec:conj}, the symbols being again in their socle degrees.
This uses Lemma~\ref{lem:mu} (the family stays in $\VZ$ and the setting holds along it), (A), Lemma~\ref{lem:det},
Lemma~\ref{lem:PF}(1), (4), (G) and the derivative rule from (R), and no hypothesis on $r$. Now let $P\in S_d$, put
$Q':=PL_a\in S_{\sigma'}$ and sum over $a$: since $Q(Lw)=LQ(w)$ and $\mathfrak D_{w}(LP')=\mathfrak D_{Lw}(P')$, the sum
is \eqref{eq:A2} for $\dot v$.

(ii) Here $Q(\dot v)=m\,f_ih_j$, $\theta=m\operatorname{adj}(M)_{ji}$ and $\mathfrak D_{\dot v}=0$. From
$g=Mh$, $h_j\det M=\sum_{i'}\operatorname{adj}(M)_{ji'}g_{i'}$. Let $\Phi^1_{(\nu)}$ be $\Phi^1=(f,p;g,s)$ with the members
of its $\nu$-th pair squared. Modulo $\Phi^1_{(\nu)}$ the product $f_ig_{i'}$ vanishes unless $i=i'$ and $\nu$ is
the pair $(f_i,g_i)$, so $f_ih_j\det M\equiv\operatorname{adj}(M)_{ji}f_ig_i$ for that $\nu$ and $\equiv0$ for the
others; by (R), for $B\in S_{d+\tau}$,
$\sum_\nu\Res_S[m f_ih_jB\det M/\Phi^1_{(\nu)}]=\Res_S[mB\operatorname{adj}(M)_{ji}/\Phi^1]$. By Lemma~\ref{lem:PF}(3) and
(1) for $Y_1$ this says $(k+1)\varphi_1(Q(\dot v)B\det M)=\pi_1(B\theta)$, the first identity \eqref{eq:A2} for $B=PN$.
The second follows in the same way from $q=M^\T f$, $f_i\det M=\sum_{j'}\operatorname{adj}(M)_{j'i}q_{j'}$ and
$\Phi^2=(h,p;q,s)$. (This step is valid for every entry of $\dot M$ and every $r$.)

Steps~B and~C of the proof of Theorem~\ref{thm:FA2}(3) use only \eqref{eq:A2}, the key lemma, Lemma~\ref{lem:det},
the relation $A_0=-\lz B_0$ (Theorem~1, through Theorem~\ref{thm:FA3}(1)) and the decomposition $N\det M=H_1+H_2$
of the proof of Theorem~\ref{thm:FA3}(1); they give the formula for each pair $N,N'$ with classes in $E'$. (The
last sentence of Step~C, on representatives, uses the range; outside the range it is replaced by the
parenthesis in the proof of Theorem~\ref{thm:FA3}(2), which does not involve the sign of $r$.) In Theorem~\ref{thm:FA3}(3) the sign of $r$ entered only through the formula.
\end{proof}

\run{Remark on $r<0$.} For $r<0$ one has $\tau\le d<\sigma'=\sigma-\tau$, where $\sigma$ is the socle degree of $\cA$; the Hilbert function
of the complete intersection $\cA$ is symmetric about $\sigma/2$ and unimodal, so $\dim\cA_\tau\le\dim\cA_d$, and
$E'=0$ whenever multiplication by $\det M$ is injective on $\cA_\tau$. Referee~24 expected this for general data. It holds at random data modulo $32003$
and $30011$ in all $19$ degree types with $r<0$ that were tried ($k\le5$, $c\le2$; coordinating thread, after
version~20), hence for general data of these types. So, at least in these types, for $r<0$ the
formula has content only at special $X$, such as the Fermat point of the Checks below (two planes with different pairings of the coordinates:
$\dim E'=2$; for two planes with the same pairing referee~25 finds $\dim E'=3$). No
corollary of this paper uses $r<0$: all rows of Table~\ref{tab:lin} and Corollary~\ref{cor:eight} have $r\ge1$.


\run{Checks of Proposition~\allr.} (a) At a Fermat point of type $(d,k,c)=(7,2,1)$ (two planes in $\PP^5$ meeting in a line, with the pairings
$(01)(23)(45)$ and $(01)(24)(35)$ of the coordinates: $r=-1$, $\tau=2$), by
the pole-order reduction of Remark~\ref{rem:r2}(d), modulo $30059$: the law holds, and the formula agrees with
the reduction for $40$ random virtual deformations of weight $1$ ($1440$ comparisons, $288$ non-zero); likewise
for $(6,2,1)$, where $r=0$ ($1440$ comparisons, $324$ non-zero). As in Remark~\ref{rem:pw}(iii), at such points
this tests the algebra that leads from Lemma~\ref{lem:PF} to the formula, and the conventions, not
Lemma~\ref{lem:PF}(3) itself. (b) Given Lemma~\ref{lem:PF}(1), (3) and Lemma~\ref{lem:det}, the identities \eqref{eq:A2} amount to two
identities between residue symbols of forms,
\[
\sum_\nu\Res_S\Bigl[\frac{Q(\dot v)PN\det M}{\Phi^i_{(\nu)}}\Bigr]=\Res_S\Bigl[\frac{PN\theta}{\Phi^i}\Bigr]\pm_i\,\mathfrak D_{\dot v}(PN),
\qquad \pm_1=+1,\ \pm_2=-\lz,
\]
where $\Phi^i_{(\nu)}$ is $\Phi^i$ with the members of its $\nu$-th pair squared. These are pure algebra and hold whether or not Lemma~\ref{lem:PF}(3) does (for tuples with $\dot M=0$, and for
$r\ge0$, they are the derivative of the identity of Lemma~\ref{lem:det}; for the matrix part with $r<0$ they are
the Cramer identity of step (ii)). So testing them checks the bookkeeping of the proposition only: degrees,
signs, and the entries $\dot M_{ij}$ at positions where $M_{ij}$ has negative formal degree. At random data (all sequences regular; $X$ not tested for smoothness, which these identities do
not need): modulo $32003$ and $30011$, with four random pairs $(P,N)$ per run, for
$r<0$ in $43$ runs of $22$ degree types ($k=2,3,4$; $c=1,2$; $-4\le r\le-1$; in five of the types an entry of $M$ has
negative formal degree) and for $r\ge0$ in $22$ runs of $11$ types: both identities hold in all $260$ instances,
with non-zero values, and fail if the $\mathfrak D$-term is dropped. Six further types with $r<0$, among them the
three with $c=3$, did not finish within the time limit.
Referee~24 tested the same identities with its own program in $31$ instances of $8$ types, $5$ of them with
$r<0$. (c) Not tested: the law itself at a random point with $r<0$ (one attempt with the implementation of
Remark~\ref{rem:r2}(e) did not finish).

\begin{corollary}[cubics outside the range]\label{cor:eight}
\begin{enumerate}
\item Let $X\subset\PP^9$ be a smooth cubic containing two $4$-planes $\Pi_1,\Pi_2$ meeting in a plane
($k=4$, $c=2$, $\tau=1=r$, $\lz=-1$). Then $\NL([\Pi_1]+\lambda[\Pi_2])$ is non-reduced, hence not smooth, at $X$
for every $\lambda\notin\{0,-1\}$, and smooth for $\lambda=-1$. For general such $X$,
$\NL([\Pi_1]+\lambda[\Pi_2])_{\mathrm{red}}=\VZ$ for $\lambda\notin\{0,-1\}$.
\item The same holds for a smooth cubic $X\subset\PP^7$ containing two $3$-planes meeting in a plane
($k=3$, $c=1$, $\tau=2=r$), with $\lz=+1$ in place of $-1$.
\item The same holds for a smooth cubic $X\subset\PP^9$ containing two $4$-planes meeting in a $3$-plane
($k=4$, $c=1$, $\tau=2$, $r=1$), with $\lz=+1$.
\end{enumerate}
\end{corollary}
\begin{proof}
In (1) and (2), $\cA=\CC[y_1,y_2,y_3]/(\bar s_1,\bar s_2,\bar s_3)$ (restriction to $\Pi_1\cap\Pi_2$), with Hilbert
function $1,3,3,1$, $\dim\cA_d=1$, and (S) holds if and only if $\det M\notin(f,h,p,s)$; then $\dim E'=2$.
(1) $\et=3$, $\sigma'=2$; the virtual deformations of weight $-1$ include $\dot M=E_{ij}$, with $Q=f_ih_j$ and
$\mathcal Q_0(N,N')(f_ih_j)=\frac\kappa5\varrho(NN'\operatorname{adj}(M)_{ji})$, and $\dot s_l$ a linear form.
(2) $\et=\dim\cA_2=3$, $\sigma'=1$, $K=\langle p_1,p_2,p_3\rangle$ and
$\mathcal Q_0(N,N')(p_l)=-\frac\kappa4\Res_S[NN'M_{00}/(f,p;h,\dots,s_l^2,\dots)]$. Exact computations over $\QQ$ on
restricted data (files \texttt{fa3\_check/}; reproduced with independent code by referee~19, files
\texttt{referee19/}): in six, respectively five, random instances the quadrics span all binary quadrics on
$E'$, so they have no common zero. Smooth instances with (S) and $\mathcal Q_0\ne0$ exist (Fermat cubics with
planes from two different matchings: referee~19). (3) $\cA=\CC[y_1,\dots,y_4]/(\bar s_1,\dots,\bar s_4)$, with Hilbert
function $1,4,6,4,1$; $\et=6$, $\dim\cA_3=4$, $\det M=M_{00}$ is linear, and (S) says $\bar M_{00}\cA_2=\cA_3$; then
$\dim E'=2$. Here $\dot M_{00}\in\CC$ gives $\mathcal Q_0(N,N')(f_0h_0)=\frac\kappa5\varrho(NN')$ and $\dot s_l$ a linear
form $\ell$ gives $-\frac\kappa5\Res_S[NN'M_{00}\ell/(\dots,s_l^2,\dots)]$; in three random instances (file
\texttt{fa3\_check/eightfold\_3plane.py}) (S) holds and these quadrics span all
binary quadrics on $E'$ (reproduced by referee~20 with its own exact code, at these and two further
instances, with all $17$ quadrics of the two kinds); smooth $X$ of this kind exist (the Fermat cubic; (S) fails
at the natural pairs of planes there, which does not matter for the closure argument). Apply Theorem~\ref{thm:FA3}(3), (4), the closedness of the
non-reduced locus, and Theorem~2(4) at $\lz$.
\end{proof}

\begin{proposition}[referee 19]\label{prop:k5}
For cubics containing two $k$-planes that meet in codimension two ($c=2$, $\tau=1$) with $k\ge5$, the map
$\cdot\det M\colon\cA_1\to\cA_3$ is injective for general data. Hence at a general smooth such $X$ there is no
excess for $\lambda\notin\{0,-1\}$: $\NL([\Pi_1]+\lambda[\Pi_2])=\VZ$ is smooth for every $\lambda\ne0,-1$, and for
$\lambda=-1$ it is the larger smooth locus of Theorem~2(4).
\end{proposition}
\begin{proof}
$\cA=\CC[y_1,\dots,y_{k-1}]/(\bar s)$ with $k-1$ quadrics, and $\bar m=\det M|_{\Pi_1\cap\Pi_2}$ is a quadric of
rank $\le4$. Injectivity is an open condition on the irreducible space of data; it holds for
$\bar s=(y_i^2)$, $\bar m=e_2(y_1,\dots,y_4)$ (a quadric of rank $4$, hence a $2\times2$ determinant): on
$y_1,\dots,y_4$ the map is the incidence matrix of points and triples of $\{1,2,3,4\}$, of determinant $-3$,
and $y_a\mapsto y_ae_2$ for $a>4$. Now apply Theorem~2(1), (4), (5).
\end{proof}

\run{Remarks.} (v) Theorem~\ref{thm:FA3} does not control the forms $Q\in K_\lambda\smallsetminus K$, which exist
exactly when (S) fails. (vi) As digested, \cite[Conj.~6, 7]{Mo} assert, for cubics of dimension $n$ containing
two linear subspaces meeting in codimension two, that the Hodge locus is smooth exactly for the coefficients
$(1,-1)$; they are stated at the Fermat point. At every smooth $X$ this dichotomy holds for $n=6$
(Corollary~\ref{cor:six}) and $n=8$ (Corollary~\ref{cor:eight}); by Proposition~\ref{prop:k5} it
fails at a general point for $n\ge10$, where the locus is smooth for every $\lambda\ne0$ (as digested, this
agrees with \cite[Thm.~1.2, 4.9]{Kl}). At the Fermat point itself, with planes of the same matching, the
quadrics vanish and nothing is decided here. No person has compared these statements with the sources.
(vii) \emph{Planes meeting in codimension one} ($c=1$, $\lz=1$). As digested by a scout agent, \cite{Kl2} proves
that for $d=3$, $k>4$, for $d=4,5$, $k>2$, and for $d\ge6$, $k\ge2$ there is an $X$ at which the two loci coincide
as schemes for all but finitely many $\lambda$ (for all $\lambda\ne0,1$ if $d\ge6$), and remarks
\cite[Remark~4.9]{Kl2} that for $(d,k)\in\{(3,3),(3,4),(4,2),(5,2)\}$ there is always excess tangent dimension,
leaving open whether the locus is then non-reduced or of larger dimension. (Its abstract states the coincidence
for all $k>1$, without these bounds.) These four pairs, with the divisor case $(3,2)$, are exactly the pairs
with $k\ge2$ and $\dim\cA_2>\dim\cA_d$, in agreement with Theorem~2(5). If the digest is accurate, Theorem~\ref{thm:FA3} answers that question:
in all four cases the locus is non-reduced at every smooth $X$, for every $\lambda\notin\{0,1\}$, and for
general $X$ its reduction is $\VZ$. For $(3,3)$ and $(3,4)$ this is Corollary~\ref{cor:eight}(2), (3). For
$(4,2)$ and $(5,2)$ (two planes meeting in a line on a quartic or quintic fourfold; $r=2$, resp. $1$):
$\cA=\CC[y_1,y_2]/(\bar s_1,\bar s_2)$ with $\deg\bar s_l=d-1$, (S) holds for general data, $\dim E'=2$, resp. $1$,
and the quadrics $N\mapsto\varrho(N^2\dot M_{00})$ and $N\mapsto\Res_S[N^2M_{00}\dot s_l/(\dots,s_l^2,\dots)]$ span all
binary quadrics, resp. are nonzero (exact, three random instances each, restricted data: file
\texttt{fa3\_check/fourfold\_line.py}; reproduced by referee~21 with independent exact code, files
\texttt{referee21/}, and by hand at $\bar s=(y_1^3,y_2^3)$, $\bar m=y_1y_2$, resp. $\bar s=(y_1^4,y_2^4)$,
$\bar m=y_1^3+y_2^3$, $N=y_1y_2$, $\dot M_{00}=y_1y_2$; for both also modulo primes with unrestricted data, see
the Checks). These two cases were treated by direct computation in \cite{P1}.

\run{Checks} (modulo primes; file \texttt{formula\_g3/REPORT\_FG3.md}; same program module as before). Outside
the range, parts (0), (1), (2) of Theorem~\ref{thm:FA3} were tested at $63$ instances with $X$ smooth modulo $p$,
of $39$ degree types with $k\le3$, $c\le4$, $\tau\le6$, $r\in\{2,4\}$, $\dim E'\le10$, $\dim\cA_d\le8$ (nine of the
instances were built so that (S) fails): in all of them the excess of $\lambda\ne\lz$ is $\dim E'$ and
$\dim K_\lambda-\dim K$ is as in (0) (five values of $\lambda$); in the $61$ with $E'\ne0$ the law holds and the
formula is exact, with the same normalisation constant as in the range, on $2284$ test triples ($1687$ nonzero).
For generic data (S) held in every type tried, and the quadrics had no common zero in $46$ of $52$ runs (in the
other six, $\dim K=1<\dim E'$); for the type of Corollary~\ref{cor:eight}(2) they span all binary quadrics, as
in the exact computation. For $Q\in K_\lambda\smallsetminus K$ (instances where (S) fails) the second-order term
does \emph{not} have $T\VZ$ in its kernel, so Remark~(v) is a real restriction. One further run (coordinating
thread; type $(d,k,c)=(5,2,1)$, the only one with $r=1$; $X$ smooth modulo $p$): (S) holds, $\dim E'=1$, the law
holds and the formula is exact on $139$ test triples, all nonzero. Not tested: $r=0$, $k=4$
(so Corollary~\ref{cor:eight}(1) and (3) rest on the proof and on the exact computations with restricted data), and
parts (3), (4) beyond the first-order dimensions.


\subsection{Two linear subspaces: the complete list at a general point}\label{sec:lin}

Let $\Pi_1,\Pi_2\subset X$ be $k$-planes meeting in dimension $k-c$ ($1\le c\le k+1$) on a hypersurface of degree
$d\ge3$. Here all $f_i,h_j,p_l$ are linear, $\tau=2c-(c-1)d$, $\lz=(-1)^{c-1}$, and
$\cA=\CC[y_1,\dots,y_{k+1-c}]/(\bar s)$ with $k+1-c$ forms of degree $d-1$ (the restrictions of the $s_l$ to
$\Pi_1\cap\Pi_2$). We assume $kd-2k-2\ge1$. In the remaining cases part (3) below fails: for cubic surfaces every $\NL(\delta)$ is
the whole space; for cubic fourfolds and quartic surfaces $V_1,V_2$ are transversal smooth divisors, $K=0$, and
$\NL(\delta_\lambda)$ is a smooth divisor strictly containing $\VZ$ for every $\lambda\ne0$. Smooth $X$ exist for all $(d,k,c)$ (Fermat).

\begin{theorem}[written by the coordinating thread; read by referee 21]\label{thm:lin}
Let $d\ge3$, $k\ge1$, $1\le c\le k+1$ with $kd-2k-2\ge1$.
\begin{enumerate}
\item If $\tau<0$ (that is: $c\ge4$; or $c=3$, $d\ge4$; or $c=2$, $d\ge5$), then
$\NL([\Pi_1]+\lambda[\Pi_2])=\VZ$ for every $\lambda\ne0$, at every smooth $X$.
\item If $\tau\ge0$, then $\NL([\Pi_1]+\lz[\Pi_2])$ is smooth and strictly larger than $\VZ$, at every smooth $X$.
\item Let $\tau\ge0$ and $\lambda\notin\{0,\lz\}$, and let $X$ be general among the smooth hypersurfaces
containing two such planes. If $(d,k,c)$ is not in the list
\[
\mathcal L=\{(3,3,1),(3,3,2),(3,3,3),(3,4,1),(3,4,2),(3,4,3),(4,2,1),(4,2,2),(5,2,1)\}\cup\{(d,1,1):d\ge5\},
\]
then $\NL([\Pi_1]+\lambda[\Pi_2])=\VZ$ as schemes (smooth). If $(d,k,c)\in\mathcal L$, then
$\NL([\Pi_1]+\lambda[\Pi_2])$ is $\VZ$ with a non-reduced structure: its reduction is $\VZ$, and its Zariski tangent
space has dimension $\dim\VZ+e$ with $e$ as in Table~\ref{tab:lin}. For $(d,k,c)\in\mathcal L$ the locus is
non-reduced at \emph{every} smooth $X$ containing two such planes.
\end{enumerate}
\end{theorem}

\begin{table}[ht]
\centering\small
\begin{tabular}{@{}llrrrl@{}}
\toprule
$(d,k,c)$ & the two planes & $\tau$ & $\lz$ & $e$ & source of the second-order statement\\
\midrule
$(3,3,1)$ & $3$-planes in $\PP^7$ meeting in a plane & $2$ & $1$ & $2$ & Cor.~\ref{cor:eight}(2)\\
$(3,3,2)$ & $3$-planes meeting in a line & $1$ & $-1$ & $2$ & Cor.~\ref{cor:six}; \cite{P1}\\
$(3,3,3)$ & $3$-planes meeting in a point & $0$ & $1$ & $1$ & Cor.~\ref{cor:Anz}; \cite{P1}\\
$(3,4,1)$ & $4$-planes in $\PP^9$ meeting in a $3$-plane & $2$ & $1$ & $2$ & Cor.~\ref{cor:eight}(3)\\
$(3,4,2)$ & $4$-planes meeting in a plane & $1$ & $-1$ & $2$ & Cor.~\ref{cor:eight}(1)\\
$(3,4,3)$ & $4$-planes meeting in a line & $0$ & $1$ & $1$ & Cor.~\ref{cor:Anz}; \cite{P1}\\
$(4,2,1)$ & planes in $\PP^5$ meeting in a line & $2$ & $1$ & $2$ & Remark~(vii); \cite{P1}\\
$(4,2,2)$ & planes meeting in a point & $0$ & $-1$ & $1$ & Cor.~\ref{cor:Anz}; \cite{P1}\\
$(5,2,1)$ & planes in $\PP^5$ meeting in a line & $2$ & $1$ & $1$ & Remark~(vii); \cite{P1}\\
$(d,1,1)$, $d\ge5$ & two coplanar lines & $2$ & $1$ & $1$ & Cor.~\ref{cor:surf}; cf.\ Maclean, Dan\\
\bottomrule
\end{tabular}
\caption{The triples $(d,k,c)$ for which the Hodge locus of $[\Pi_1]+\lambda[\Pi_2]$, $\lambda\notin\{0,\lz\}$, differs
from $\VZ$ at a general point: there it is $\VZ$ with a non-reduced structure and $e$ extra tangent
directions.}\label{tab:lin}
\end{table}

\begin{proof}
(1), (2) are Theorem~2(3), (4); in (2) the locus is strictly larger because $0\le\tau\le\sigma=\tau+(kd-2k-2)$,
so $\et=\dim\cA_\tau\ge1$. (3) By Theorem~2(5) the excess of $\lambda$ is
$\dim E'=\dim\ker(\cdot\bar m\colon\cA_\tau\to\cA_d)$, $\bar m:=\det M|_{\Pi_1\cap\Pi_2}$, and if $E'=0$ then
$T\NL(\delta_\lambda)=T\VZ$, so $\NL(\delta_\lambda)=\VZ$. Injectivity of $\cdot\bar m$ is an open condition on the
irreducible space of data, so one instance suffices, and $\tau\ge0$ leaves three families.

$c=1$ ($\tau=2$, $\bar m$ an arbitrary form of degree $d-2$, $\cA$ in $k$ variables, socle degree $k(d-2)$): take
$\bar s=(y_i^{d-1})$ and $\bar m=L^{d-2}$, $L:=\sum y_i$. By the strong Lefschetz property of monomial complete
intersections in characteristic $0$ (\cite{St,Wa}; not opened), $L^{\sigma-2i}\colon\cA_i\to\cA_{\sigma-i}$ is
bijective for $i\le\sigma/2$; hence $L^j\colon\cA_i\to\cA_{i+j}$ has maximal rank for all $i,j$ (if
$i+j\le\sigma-i$ it is a factor of an injective map; otherwise it is dual to
$L^j\colon\cA_{\sigma-i-j}\to\cA_{\sigma-i}$, which is injective by the first case). So $\cdot\bar m\colon\cA_2\to\cA_d$
is injective if and only if $\dim\cA_2\le\dim\cA_d$. If $d\le k(d-2)-2$ this holds by symmetry and unimodality of
the Hilbert function. Otherwise $(k-1)(d-2)\le3$, i.e. $k=1$ or
$(d,k)\in\{(3,2),(3,3),(3,4),(4,2),(5,2)\}$, and $\dim\cA_2>\dim\cA_d$ by inspection ($k=1$: $\cA_d=0\ne\cA_2$; the
Hilbert functions are $1,2,1$; $1,3,3,1$; $1,4,6,4,1$; $1,2,3,2,1$; $1,2,3,4,3,2,1$); $(3,2)$ is excluded by
$kd-2k-2\ge1$.

$c=2$, $d=3$ ($\tau=1$): injective for general data if and only if $k\ge5$ (Proposition~\ref{prop:k5} and the
dimension count); $k=2$ is excluded.
$c=2$, $d=4$ ($\tau=0$, $\cA$ in $k-1$ variables with cubic relations, $\bar m$ the determinant of a $2\times2$ matrix
of quadrics): $E'=0$ if and only if $\bar m\ne0$ in $\cA_4$, which holds for general data if and only if
$\cA_4\ne0$ (products of two quadrics span the quartics), i.e. $2(k-1)\ge4$; so only $k=2$ remains ($k=1$ is
excluded).
$c=3$, $d=3$ ($\tau=0$, $\cA$ in $k-2$ variables with quadratic relations, $\bar m$ the determinant of a $3\times3$
matrix of linear forms): $E'=0$ for general data if and only if $\cA_3\ne0$, i.e. $k-2\ge3$; so $k=3,4$ remain
($k=2$ is excluded).

For the triples of $\mathcal L$ the values of $e=\dim E'$ at a general point follow from the Hilbert functions
(and from (S), which holds for general data: trivially in the range; for $c=1$ by the Lefschetz instance above,
since maximal rank is surjective when $\dim\cA_2\ge\dim\cA_d$; for $(3,4,2)$ because (S) there says
$\bar m\notin(\bar s)_2$), and the second-order statement is proved in the
places listed in Table~\ref{tab:lin}: for $\tau=0$ the type is in the range, $\et=1$ and $A\ne0$ for general data
(Corollary~\ref{cor:Anz}); for coplanar lines Corollary~\ref{cor:surf}; for $(3,3,2)$ Corollary~\ref{cor:six}; for
the others Theorem~\ref{thm:FA3}(3), (4) with the exact computations of Corollary~\ref{cor:eight} and
Remark~(vii) or, without a computer, with the monomial instances of Remark~(xi), at which $(f,h,p,s)$ is
regular, (S) holds and the quadrics have no common zero.
Non-reducedness at every smooth $X$ follows because the non-reduced locus is closed and $\mu$ is open onto
$\VZ$.
\end{proof}

\run{Remarks.} (viii) Parts (1), (2) and the first half of (3) use the period formula only through
Lemma~\ref{lem:PF}(1), (2) ($\kappa\ne0$), not the constancy of $\kappa$; the non-reduced half of
(3) uses Lemma~\ref{lem:PF}(4). Neither uses (CG). The first half of (3) also uses the Lefschetz property quoted in the
proof (checked directly for $d\le10$, $k\le9$: Remark~\ref{rem:r2}(b)). (ix) As digested, the first-order thresholds
agree with \cite{Kl} (Theorems~1.2, 1.3, 4.9) and \cite{Kl2} (Theorem~1.1, Remark~4.9), where equality of the two
loci is proved at some $X$ for all but finitely many $\lambda$, and with Villaflor Loyola's theorem for
$\tau<0$ at the Fermat point; for coplanar lines the non-reduced statement is, in the cases they treat and as
digested, due to Maclean and Dan; the cases $(3,3,3)$, $(3,4,3)$ are those of \cite[Conj.~1]{Mo} and of the
alternative left open in \cite[Thm.~1.3]{Kl} (as digested), and were treated in \cite{P1}; the four cases with
$c=1$, $k\ge2$ are those of \cite[Remark~4.9]{Kl2} (Remark~(vii)). We have not found the
non-reduced statements for the other triples of $\mathcal L$, nor the list $\mathcal L$ as a whole, in the
sources read (machine digests). (x) At special points the picture changes: $E'$ can be larger, and the
quadrics can vanish (the Fermat hypersurface with planes of the same matching). (xi) \emph{Instances that can
be checked by hand} (referee~21), for the five rows of Table~\ref{tab:lin} outside the range; in each,
$\bar s$ is monomial, so the residues are coefficients:
$(4,2,1)$: $\bar s=(y_1^3,y_2^3)$, $\bar m=y_1y_2$, $E'=\langle y_1^2,y_2^2\rangle$, quadrics $2n_1n_3,n_1^2,n_3^2$;
$(5,2,1)$: $\bar s=(y_1^4,y_2^4)$, $\bar m=y_1^3+y_2^3$, $E'=\langle y_1y_2\rangle$, value $1$ at $\dot M_{00}=y_1y_2$;
$(3,3,1)$: $\bar s=(y_i^2)$, $\bar m=\sum a_iy_i$, $E'=\{a_3n_{12}+a_2n_{13}+a_1n_{23}=0\}$, quadrics
$a_1n_{12}n_{13}$, $a_2n_{12}n_{23}$, $a_3n_{13}n_{23}$ (no common zero iff $a_1a_2a_3\ne0$);
$(3,4,1)$: $\bar s=(y_i^2)$, $\bar m=\sum y_i$, $E'=\{n_{12}=n_{34}=\alpha,\ n_{13}=n_{24}=\beta,\ n_{14}=n_{23}=\gamma,\
\alpha+\beta+\gamma=0\}$, quadrics $\alpha\beta,\alpha\gamma,\beta\gamma,\alpha^2+\beta^2+\gamma^2$;
$(3,4,2)$: $\bar s=(y_i^2)$, $\bar M=\bigl(\begin{smallmatrix}y_1&y_2\\y_3&y_1+y_2+y_3\end{smallmatrix}\bigr)$,
$E'=\{n_1=n_2+n_3\}$, quadrics $n_1n_2,n_1n_3,n_2n_3$. The polynomial algebra of these five instances (the
equations of $E'$, the quadrics as coefficients of the stated products, and identities showing that the
squares of the coordinates on $E'$ lie in the ideal of the quadrics) is certified by the Lean~4 kernel (file
\texttt{LinearListCert.lean}, $63$ theorems, generated by an agent and re-run by the coordinating thread); it
certifies nothing about residues, Hodge loci or the passage from these instances to general data.


\section{Computations}\label{sec:comp}

All computations in this section are modulo primes at random instances; they are evidence, not proofs. The main
program (primes $2^{31}-1$, $2147483629$, $10^9+7$; $X$ certified smooth modulo the prime) computes the period
polynomials of $Y_1,Y_2$, the tangent spaces, the space $K$ and the quadratic term, in the Jacobian ring.

\begin{table}[h]\centering\small
\begin{tabular}{llll>{\raggedright\arraybackslash}p{0.2\textwidth}>{\raggedright\arraybackslash}p{0.2\textwidth}}
\toprule
$(k,d,c)$; $\deg f$; $\deg h$; $\deg p$ & $\tau$ & range & $\et$ & excess found & quadratic map on the excess
directions vanishes at\\
\midrule
$(1,6,1)$; $(1)$; $(1)$; $(2)$ & $2$ & yes & $2$ & $2$ & $+1$ only\\
$(1,6,1)$; $(1)$; $(1)$; $(3)$ & $2$ & yes & $3$ & $3$ & $+1$ only\\
$(1,5,1)$; $(1)$; $(1)$; $(2)$ & $2$ & yes & $2$ & $2$ & ($K=0$)\\
$(1,5,1)$; $(1)$; $(1)$; $(3)$ & $2$ & yes & $2$ & $2$ & ($K=0$)\\
$(1,5,1)$; $(1)$; $(1)$; $(1)$ & $2$ & yes & $1$ & $1$ & $+1$ only\\
$(2,3,2)$; $(1,1)$; $(1,1)$; $(1)$ & $1$ & yes & $1$ & $1$ & ($K=0$)\\
$(2,3,1)$; $(1)$; $(1)$; $(1,1)$ & $2$ & yes & $1$ & $1$ & ($K=0$)\\
$(2,4,3)$; $(1,1,1)$; $(1,2,3)$; -- & $1$ & yes & $2$ & $2$ & $+1$ only\\
$(2,4,3)$; $(1,1,2)$; $(1,2,2)$; -- & $1$ & yes & $3$ & $3$ & $+1$ only\\
$(2,4,2)$; $(1,1)$; $(1,2)$; $(1)$ & $1$ & yes & $2$ & $2$ & $-1$ only\\
$(2,4,1)$; $(1)$; $(1)$; $(1,1)$ & $2$ & no & $3$ & $3$ at $\lambda=+1$, $2$ at the six others & $+1$\\
$(2,4,3)$; $(1,1,1)$; $(1,1,2)$; -- & $-1$ & -- & $0$ & $0$ & --\\
\bottomrule
\end{tabular}
\caption{Theorem 2 against runs made before it was formulated ($\tau\ne0$). Last column: among the tested
$\lambda\in\{1,-1,2,3,\frac12,-2,\frac53\}$, the values at which the quadratic map on the excess directions
vanishes. In the range the excess is predicted to be $\et$; outside it, $\et$ at $\lz$ and
$\dim(0:_\cA\det M)_\tau$ elsewhere (in row 11: $3-1=2$).}\label{tab:tau}
\end{table}

\begin{table}[h]\centering\small
\begin{tabular}{l>{\raggedright\arraybackslash}p{0.55\textwidth}>{\raggedright\arraybackslash}p{0.22\textwidth}}
\toprule
$(k,d)$ & types with $\tau=0$, $K\ne0$ ($\deg f+\deg h$; shared $p$) & found\\
\midrule
$(1,6),(1,7),(1,8)$ & $(1,2)+(1,2)$; $(1,2)+(2,2)$; $(1,1)+(3,3)$ \quad($c=2$) & $\lz=-1$, $A\ne0$\\
$(2,4)$ & $(1,1)+(1,1)$ with $p$ of degree $1$, $2$, $3$ \quad($c=2$) & $\lz=-1$, $A\ne0$\\
$(2,4)$ & $(1,1,1)+(1,1,3)$; $(1,1,1)+(1,2,2)$; $(1,1,2)+(1,1,2)$ \quad($c=3$) & $\lz=+1$, $A\ne0$\\
$(2,5)$ & $(1,1)+(1,2)$ with $p$ of degree $1$, $2$ \quad($c=2$) & $\lz=-1$, $A\ne0$\\
$(2,5)$ & $(1,2,2)+(1,2,2)$; $(1,1,1)+(1,3,3)$; $(1,1,3)+(1,1,3)$; $(1,1,1)+(2,2,3)$;
$(1,1,2)+(2,2,2)$; $(1,1,2)+(1,1,4)$ \quad($c=3$) & $\lz=+1$, $A\ne0$\\
$(3,3)$ & $(1,1,1)+(1,1,1)$, $p$ linear \quad($c=3$) & $\lz=+1$, $A\ne0$\\
$(3,3)$ & $(1,1,1,1)+(1,1,1,2)$ \quad($c=4$) & $\lz=-1$, $A\ne0$\\
$(4,3)$ & $(1,1,1,1)+(1,1,1,2)$, $p$ linear \quad($c=4$) & first order only: $\lz=-1$\\
$(4,3)$ & $(1,1,1,1,1)+(1,1,1,2,2)$; $(1,1,1,1,2)+(1,1,1,1,2)$ \quad($c=5$) & first order only: $\lz=+1$\\
\bottomrule
\end{tabular}
\caption{The equality case: in every run the excess is $1$ and $K$ is independent of $\lambda$; except in the
last two rows the quadratic term vanishes exactly at $\lz=(-1)^{c-1}$ among the tested values, and the law
\eqref{eq:law} holds exactly modulo the prime with $A\neq0$. ``First order only'': $\det M\in T\NL(\delta_\lambda)$
exactly at $\lz$. Not exhaustive: e.g. for $(k,d,c)=(2,5,3)$ the types $(1,1,1)+(1,2,4)$, $(1,1,2)+(1,2,3)$,
$(1,1,3)+(1,2,2)$, and for $(2,5,2)$ $p$ of degree $3$, $4$, were not run. For $d=5$: one run per type, five
values of $\lambda$, and a column subset of the linear system.}\label{tab:types}
\end{table}

\run{Checks for $c=4,5$ and of $R$} (one verifier agent and three referees, each with its own code).
(i) \emph{The family $W_t$} (modulo $2^{31}-1$). For $(k,d,c)=(3,3,4)$, $(4,3,4)$ with $p$ linear, and the two
types with $(k,d,c)=(4,3,5)$, and for both choices of $h_n$ when the degrees differ: $F-t\det M\in I(t)_d$ (and,
as a control on the sign convention, $F+t\det M\notin I(t)_d$); the Hilbert function of $S/I(t)$ for $t=1,2$
and a random $t$ (for $k=4$ with $h_n$ linear: $t=1$ only) equals that of $S/I(0)=S/(I(Y_1)\cap I(\Theta))$ in
all degrees computed ($\le8$, $7$, $6$ for $8$, $9$, $10$ variables), with leading differences
$\deg Y_1+\deg\Theta=17$ and $5$ (resp. $34$ and $10$); $X$ is smooth, the tangent spaces have the codimensions
of Theorem~2, and $\det M\in T\NL(\delta_\lambda)$ exactly for $\lambda=(-1)^{c-1}$ among seven tested values. The
second-order law was run only for $(3,3,4)$ (two instances: it holds with $\lz=-1$, $A\ne0$); for $k=4$ it was
too large, so for $c=5$ the value $\lz=+1$ is supported at first order and by the Hilbert functions only.
(ii) \emph{Outside the range} (modulo $2^{31}-1$): $(k,d,c)=(3,4,4)$, $f$ linear, $h$ quadrics, $\tau=0$: $X$
smooth, $TV_1+TV_2$ of codimension $1$ in $S_d$, $T\NL(\delta_\lambda)=TV_1\cap TV_2$ for the six tested
$\lambda\neq-1$ and one dimension more at $\lambda=-1$, with $\det M$ in it: as Theorems~1 and~2(4), (5) predict
(for general $M$, $\det M\notin(f,h,p,s)_d$).
(iii) \emph{The residual $R$:} in five types including $c=4$ (modulo $2^{31}-1$), $\dim R=k$, $\deg R$ is as
predicted by Theorem~1, and $(p,q):\det M$ has the Hilbert function of $Y_1$; the two colon formulas of
Lemma~\ref{lem:NW} hold in all degrees tested ($\le3$, $4$ or $5$) in eight types, among them one with $c=1$,
one with $\tau=1$, one with $\tau=-1$ and one with a constant column in $M$ (modulo $10^9+7$; smoothness of $X$ not tested there). By hand
(referee~12, re-computed by referee~13): for skew lines on a quartic surface $\deg R=8$, $R^2=12$, $p_a(R)=7$;
for a line and a disjoint conic on a quintic surface $11$, $17$, $15$; for two disjoint planes on a cubic
fourfold $\deg R=7$, $R^2=21$. (iv) \emph{A caution on $W_t$} (referee~11): for skew lines on a quartic with
$M_{00},M_{10}$ sharing a linear factor modulo $h_1$, $S/I(t)$ has Hilbert polynomial $4m$ and its saturation
with respect to $h_n$ has $3m+1$: $W_t$ is a cubic plus a fixed line. (v) The hypothesis ``all degrees $<d$''
cannot be dropped: $(1,1,1,1,1)+(1,1,1,1,3)$, $d=3$, and $(1,1,1)+(1,1,5)$, $d=5$, have $\tau=0$ but admit no
smooth $X$. Inside the setting a smooth $X$ requires $\det M\not\equiv0$, hence a permutation $\pi$ with
$d_i+e_{\pi(i)}\le d$ for all $i$.

\begin{remark}[deduction steps checked by Lean; added in version 17, not refereed]\label{rem:lean2}
Six further Lean files ((a)--(c) added with version~17, (d) with version~20, (e), (f) with version~21) contain theorems
with hypotheses, for arbitrary data, rather than identities between explicit polynomials.

(a) \texttt{SecondOrderLaw.lean} (Lean core, no library). Let $R$ be a commutative ring, $\varphi_1,\varphi_2$
symmetric bilinear forms with values in $R$, $T_1,T_2$ subsets, $\lz^2=1$, and $D,D',G_2,G_2'$ elements with
$G_2,G_2'\in T_1$ and $G_1:=D+\lz G_2\in T_2$, $G_1':=D'+\lz G_2'\in T_2$. If $\varphi_i$ vanishes on $T_i\times T_i$
$(i=1,2)$ and $\varphi_1(D,D')+\lz\varphi_2(D,D')=0$, then for $\lambda\mu=1$, with $G_\lambda:=G_1-\mu G_2$:
\[
(\varphi_1+\lambda\varphi_2)(G_\lambda,G_\lambda')=(1-\lz\mu)\bigl[\varphi_1(G_1,G_1')-\varphi_2(G_1,G_2')-\varphi_2(G_2,G_1')\bigr].
\]
With $\varphi_i(G,G')=\varphi_i(QGG')$, $D=N\det M$ and $\mu=1/\lambda$ this is the algebra in the proofs of
Theorem~3, Theorem~\ref{thm:FA2}(2) and Theorem~\ref{thm:FA3}(1); its hypotheses are the key lemma and the
consequence $A_0+B_0/\lz=0$ of Theorem~1. The file also contains Steps~B and~C of the proof of
Theorem~\ref{thm:FA2}(3) (hypotheses: the three evaluations of Step~C; conclusion:
$(k+1)A_0=\kappa[\varrho-\mathfrak D]$), the displayed identity for $z\in T_1\cap T_2$ in the proof of
Theorem~\ref{thm:FA3}(1), and a model over $\mathbb{Z}$ in which all hypotheses hold and $A_0=1$.

(b) \texttt{SecondOrderLawModule.lean} (with Mathlib): the same statements for modules, bilinear maps (symmetry is
not needed) and submodules; that $G_\lambda\in\ker(u_1+\lambda u_2)$ if $u_1(G_2)=0$, $u_2(G_1)=0$ and
$u_1(G_1)=u_2(G_2)$; and, over a field: if $A_0(N,N)\ne0$ and $\lambda\notin\{0,\lz\}$, the second-order term at
$\lambda$ is non-zero on $G_\lambda(N)$.

(c) \texttt{SpecialCoefficient.lean} (with Mathlib). For a square matrix $M$ of any size $c$ over any commutative
ring, and $g=Mh$, $q=M^\T f$: $f^\T Mh=\sum_jq_jh_j$; $\det M\cdot h_j\in(g)$ and $\det M\cdot f_i\in(q)$ with
explicit coefficients (Cramer); the two transition matrices from $(f,h,w)$ to $(f,Mh,w)$ and to $(h,M^\T f,w)$
have determinants $\det M$ and $(-1)^c\det M$; and for every additive symbol that satisfies the transformation
law (R) (a hypothesis): the first sentence of Lemma~\ref{lem:det}, and
$\Res[m\det M/\Phi^1]+(-1)^{c+1}\Res[m\det M/\Phi^2]=0$. Here the sequences are ordered $(f,h,p,s)$ instead of $(f,p;h,s)$; that this common
reordering does not affect the signs is not formalised.

(d) \texttt{PointwiseFormula.lean} (with Mathlib; added with version~20): the algebra of
Lemma~\ref{lem:PF}(3), (4); see Remark~\ref{rem:pw}(iv).

(e) \texttt{EveryR.lean} (with Mathlib; added with version~21): step (ii) of the proof of Proposition~\allr, for a
matrix of any size over any commutative ring. Hypotheses: for the sequences $\Phi^1$ and $\Phi^2$, the symbol with
the $\nu$-th pair squared vanishes on multiples of the members of the other pairs, and
$\Res_S[x\,u_iv_i/\Phi_{(i)}]=\Res_S[x/\Phi]$ for an own pair $(u_i,v_i)$. Conclusions:
$\sum_\nu\Res_S[m f_ih_jB\det M/\Phi^\iota_{(\nu)}]=\Res_S[mB\operatorname{adj}(M)_{ji}/\Phi^\iota]$ for $\iota=1,2$; and, with
Lemma~\ref{lem:PF}(1), (3) as further hypotheses, the two identities \eqref{eq:A2} for a virtual deformation with a
single matrix entry.

(f) \texttt{SecondProof.lean} (with Mathlib; added with version~21): the second proof in Remark~\ref{rem:pw}. From
the derivative relation $(k+1)\varphi(PG)-\kappa\sum_\nu\Res_S[PG/\Phi_{(\nu)}]=-\dot\kappa\,\varrho(P)$ and the vanishing
of $\varrho$ on the ideal (hypotheses), the pointwise formula holds on the additive span of the products $PG$ with
$P$ in the ideal, and $\dot\kappa\,\varrho(P)=0$ for every $P$ with $PG$ in that span.

In all these files the geometric inputs (the key lemma, Theorem~1, the period formula, the transformation law,
regularity of the sequences) are hypotheses. The kernel certifies the deductions from them and nothing else; it
does not certify any statement about Hodge loci. Variants of the statements with a wrong sign or a missing
hypothesis are rejected. Mathlib (tag \texttt{v4.36.0}) was compiled from source for (b)--(f).
\end{remark}

\begin{remark}[checks made after version 17 by other methods; not refereed]\label{rem:r2}
Files: \texttt{CHECKS\_R2.md} and the reports of the independent implementer. All computations are modulo primes
at the instances named, with $X$ smooth where a random $X$ is used.

(a) \emph{Intersection numbers.} Under Lemma~\ref{lem:PF} and the multiplicative form of (CG), the equality
$\kappa(X,Y_1)=\kappa(X,Y_2)$ is equivalent to
\[
\rho_J(J_1J_2)=\deg Y_1\deg Y_2-d\,(Y_1\cdot Y_2),\qquad
Y_1\cdot Y_2=\deg V(f,h,p)\cdot\bigl[\eta^{k-c}\bigr]\frac{\prod_l(1+\deg p_l\,\eta)}{1+d\eta}
\]
($J_i$ the Jacobian determinants of $(f_1,g_1,\dots,p_1,s_1,\dots)$ and $(h_1,q_1,\dots,p_1,s_1,\dots)$; the
intersection number by the excess intersection formula, $0$ if there is no shared equation). This identity and
the two for the self-intersections hold at a random $X$ in each of $190$ degree types ($k\le4$, $0\le c\le k+1$,
every sign of $\tau$; Singular, modulo $32003$), and again modulo $2147483629$ in $60$ of them; with the opposite
sign the first would fail in all. So the rule $\lz=(-1)^{c-1}$ is confirmed without Theorem~1, also for $\tau<0$.
Five types with $(k,d)=(3,4)$ were too slow. The identity is insensitive to a dependence of $c(k,d)$ on $X$.

(b) \emph{The Lefschetz statement} used for $c=1$ in the proof of Theorem~\ref{thm:lin}: for
$\cA=\QQ[y_1,\dots,y_k]/(y_i^{d-1})$ the map $L^{d-2}\colon\cA_2\to\cA_d$ has maximal rank for all $3\le d\le10$,
$1\le k\le9$, by a rank computation modulo a prime (maximal rank modulo a prime implies it over $\QQ$, so these $72$
instances are proved without \cite{St,Wa}); the kernel is nonzero exactly for the pairs $(d,k)$ of the $c=1$ rows of
Table~\ref{tab:lin} and the two excluded pairs $(3,2)$, $(4,1)$, with dimension $e$.

(c) \emph{Table~\ref{tab:lin}, first order at a random point}, from
$T\NL(\delta_\lambda)=(J_F:(J_1+\lambda J_2))_d$ alone (Singular; Theorem~2 not used): the excess is the $e$ of the
table in all rows ($d=5,6,7$ in the last), and $0$ for the eight triples
$(3,3,4)$, $(3,4,4)$, $(3,4,5)$, $(4,2,3)$, $(5,2,2)$, $(5,2,3)$, $(5,1,2)$, $(6,1,2)$ outside it.

(d) \emph{Second order at Fermat points by pole-order reduction.} For $F=\sum x_i^d$ one has
$\Res(x^\beta\Omega/F^{m+1})=\frac{\beta_i-d+1}{dm}\Res(x^{\beta-de_i}\Omega/F^m)$ if $\beta_i\ge d-1$, and a
monomial form with all exponents $\le d-2$ is an eigenvector of the diagonal automorphisms, hence of pure Hodge
type. So at a Fermat point $\varphi_{[Y]}$ is determined by $\pi_{[Y]}$, i.e.\ by Lemma~\ref{lem:PF}(1), and the
second-order term of (G) can be computed without Theorems~1, \ref{thm:FA}, \ref{thm:FA2}, \ref{thm:FA3}, the key
lemma, or the constancy of $\kappa$. For two linear subspaces of the Fermat hypersurface given by equations
$x_a=\zeta x_b$ with different pairings of the coordinates ($39$ configurations, $67$ runs; all nine sporadic rows
of Table~\ref{tab:lin} and coplanar lines for $d=5,6$; primes $32003$, $50023$, $30011$, $30089$, $30013$): $T\VZ$ is in
the kernel of the second-order form; the form vanishes at $\lz$; the law of Theorems~\ref{thm:FA2}(2),
\ref{thm:FA3}(1) holds at six values of $\lambda$ in all $61$ runs with excess; and the formula of
Theorem~\ref{thm:FA2}(3) holds with its constant (in these units $\kappa=(-1)^{(k+1)(k+2)/2}$, no fitted constant)
in $25344$ comparisons, $8532$ of them with nonzero value. This includes the cubic eightfold rows $(3,4,1)$,
$(3,4,2)$, which had not been run before. The instance of Corollary~\ref{cor:six}, Remarks~(i) and~(iv) of
\S\ref{sec:FA2}, and the vanishing for two planes of the same pairing meeting in a point are reproduced. At some
Fermat configurations of $(3,4,1)$ and $(5,2,1)$ the excess is $e+1$ and $\dim K_\lambda=\dim K+1$, as
Theorem~\ref{thm:FA3}(0) requires when the rank of $\cdot\det M$ drops; at some Fermat configurations of $(3,3,1)$,
$(3,3,2)$, $(3,4,1)$, $(4,2,1)$ the quadrics have a common zero. These are special points. A Fermat point cannot show a dependence of $\kappa$ on $X$, and the comparison of
the formula there tests its algebra and its constant only: at these configurations the pointwise formula
\eqref{eq:pw} holds for formal reasons (Remark~\ref{rem:pw}(iii)).

(e) \emph{An independent implementation.} An AI agent that was given only this paper (no program, note or report)
re-implemented the definitions in Macaulay2 and tested random instances. It reported no failure in any check (raw logs kept and tallied by the coordinating thread: no failed flag; one of
its runs was repeated with identical output). First order: Theorem~2 and Lemma~\ref{lem:det} in $42$ types
($k\le4$, $1\le c\le5$, $-2\le\tau\le6$, twelve outside the range) at two primes each; the algebra of Theorem~1
(both colon formulas; dimension, degree and Cohen--Macaulayness of $R$) in $12$ types at two primes. Second
order at random points, $123$ instances of $62$ types, primes $32003$ and $30011$: the key lemma, the description
of $K$, Theorem~\ref{thm:FA3}(0), vanishing at $\lz$, the law at five or six values of $\lambda$, $T\VZ$ in the
kernel, and the formula of Theorems~\ref{thm:FA2}(3), \ref{thm:FA3}(2) as an equality (with $\kappa=1$) on $2464$
test triples, $2196$ of them nonzero, together with the two identities \eqref{eq:A2}; $35$ types are in the range
($10$ with more than one excess direction) and $19$ are outside it with $E'\ne0$ ($6$ with (S) failing). It ran the
cases that had not been run before: the cubic eightfold rows $(3,4,1)$, $(3,4,2)$ at random points ($\dim E'=2$,
(S) holds, the quadrics span all binary quadrics) and one cubic eightfold with (S) failing; the second-order term
for the $k=4$ equality types with $c=3,4,5$ ($A\ne0$; $\lz=+1$ for $c=5$); and three special types with $r=0$ and
$E'\ne0$. All ten rows of Table~\ref{tab:lin} are reproduced at second order at a random point and two primes
($e$, and no common zero), and $13$ triples outside the list, among them $(4,3,1)$, $(4,3,2)$, $(6,2,1)$,
$(6,2,2)$ and the cubic tenfolds $(3,5,c)$, $c=1,2,3$, show no excess. Special instances reproduce
Corollary~\ref{cor:surf} ($M_{00}(o)=0$: $A=0$ on a quintic, $A\ne0$ on a sextic), Corollary~\ref{cor:six}(1)
($\det M|_L$ a square: a common zero), Remark~(v) of \S\ref{sec:FA3}, all rows of Table~\ref{tab:tau} and the
codimensions quoted in the introduction. Not run: second order for $k=5$, for $r<0$, and for $r=0$ with non-linear
degrees; seven types with $(k,d)=(2,5)$ of the last table. Its evaluation of $\varphi_\delta$ rests on
\eqref{eq:pw} (now Lemma~\ref{lem:PF}(3)), which it derived itself; so these runs test the algebra of the paper given (G), (A) and
Lemma~\ref{lem:PF}(1), not those inputs. Its reading notes led to the rewording of (G), (R) and (PF) in
\S\ref{sec:facts} in version~19: how to evaluate $\varphi_\delta$; the exact coefficients in (G); the order inside
the sequences and the normalisation and direction of (R); and the remark that the sign $(-1)^{k+1}$ in the Checks of
\S\ref{sec:conj} depends on conventions (in its normalisation the ratio is $+1$ for $k\le4$).

(f) With the formula, at a random point of each row of Table~\ref{tab:lin} (two seeds): $\dim E'=e$ and the
quadrics have no common zero.
\end{remark}

\begin{remark}[history and checks of Lemma~\ref{lem:PF}(3), (4)]\label{rem:pw}
The pointwise formula \eqref{eq:pw} was found by the implementer of Remark~\ref{rem:r2}(e), who observed that
(G) defines the second-order term through $\varphi_\delta$ while the paper did not say how to evaluate
$\varphi_\delta$ outside the Jacobian ideal. Versions~18 and~19 stated it in this remark, with the theorems still
conditional on (CG). An AI referee (``referee~23'', who saw only version~19) then checked the remark line by line and with its own
computations and found no error; its corrections are applied here (the pole-order-$(k+1)$ form of (G), the
holomorphy of $\kappa$, two justifications in the proof, a clash of notation, and an overstatement in the checks,
corrected below). Since version~20 it is Lemma~\ref{lem:PF}(3), the constancy of $\kappa$ is
Lemma~\ref{lem:PF}(4), and the statements of \S\S\ref{sec:conj}--\ref{sec:lin} no longer assume (CG).
What Lemma~\ref{lem:PF}(4) does \emph{not} give: the value or the sign of $\kappa$; that $\kappa$ is the same for
different degree types; and $\kappa(X,Y_1)=\kappa(X,Y_2)$, which is still Lemma~\ref{lem:det} (Theorem~1,
$\tau\ge0$). No result needs the first two. The other cited inputs (G), (A), (SI), (R) are textbook material, but
like (CG) they were not opened by anyone here.

\emph{A second proof} (referee~23). Differentiating $\pi=\kappa\varrho_Y$ along the family of the proof gives, for
$P\in S_t$ and $G=\sum(\alpha_\nu\psi_\nu+\phi_\nu\beta_\nu)$,
$(k+1)\varphi_{[Y]}(PG)-\kappa\sum_\nu\Res_S[PG/\Phi_{(\nu)}]=-\dot\kappa\,\varrho_Y(P)$. For $P\in(\phi,\psi)_t$ the right
side vanishes, and $(\phi,\psi)_t\cdot(\phi,\psi)_d$ spans $S_{t+d}$; this is \eqref{eq:pw}, and then $\dot\kappa=0$.
It uses the pole-order-$(k+1)$ expansion of (G) only.

\emph{Checks of \eqref{eq:pw}} (modulo primes unless said otherwise; none is a proof). No test of ours compares
\eqref{eq:pw} outside the Jacobian ideal with an independently computed period: there its support is the two
proofs. What the tests do check (the assessment is that of referees~23, 24 and~25, confirmed by the coordinating
thread):
(i) \emph{Equal weights of the own-pair terms.} With \eqref{eq:pw}, the second-order term of $\delta_{\lz}$ along the
pencil $F-N\det M$, which must vanish by Theorem~1, does vanish: at all $123$ random instances of the implementer,
and for \emph{all} $Q\in S_{kd-2k-2}$ (not only $Q\in K$) in eleven runs of ten degree types with $k\le2$ (program of
referee~23, re-run by the coordinating thread; every test form lay outside the Jacobian ideal, every value was
non-zero, and the opposite sign fails). For $c=1$ the vanishing holds term by term in $\nu$, by (R), and says
nothing about \eqref{eq:pw}. For every $c$ the terms of the shared pairs $(p_l,s_l)$ cancel one by one, by (R)
(referee~25). For $c\ge2$ (eight of the eleven runs; in seven of them $c=k+1$) the $c$ terms of the own pairs
$(f_i,g_i)$, $(h_i,q_i)$ do not cancel one by one, and with unequal weights their sum is non-zero. So these runs
test, outside the Jacobian ideal, that the own-pair terms enter with equal weights (all terms if and only if
$c=k+1$), and nothing more.
(ii) \emph{Additivity.} For a generator $\phi_0=uv$ one has $[Y]=[V(u,\phi_1,\dots)]+[V(v,\phi_1,\dots)]$, and both
sides of \eqref{eq:pw} are additive (eight instances, $k\le2$; same program, re-run). This holds term by term in
$\nu$, so it is blind to weights and constants.
(iii) \emph{Constants and conventions, on the Jacobian ideal.} Agreement with Griffiths--Dwork reduction on the
Jacobian ideal is a formal identity of residue symbols; it fixes the factor $k+1$ and the equal weights there.
Exactly, by hand (referee~23; recomputed by the coordinating thread and by referee~24): for the line
$V(x_0-x_1,x_2-x_3)$ on $x_0^4-x_1^4+x_2^4-x_3^4=0$, $\Res_S[x_0^2x_2^2/(\phi;\psi)]=-\frac1{16}$ and
$\sum_\nu\Res_S[x_0^6x_2^2/\Phi_{(\nu)}]=-\frac3{64}$; the ratio $\frac34$ is $k+1$ times the factor $\frac38$ of
pole-order reduction. The agreement with the pole-order reduction of Remark~\ref{rem:r2}(d) at Fermat
points ($780$ random forms in $39$ configurations) has no content beyond this, because for a linear subspace
aligned with a pairing of the coordinates both sides vanish, by symmetry, on the monomials outside the Jacobian
ideal (referee~23; versions~18 and~19 overstated this check). Independence of the presentation
($\psi\mapsto\psi+A\phi$ with $A$ skew, $\phi\mapsto B\phi$, $\psi\mapsto(B^{-1})^\T\psi$; $21$ degree types) follows
from the algebra and tests the programs. (iv) The algebra of the proofs of (3) and (4), with the
$\epsilon^2$-coefficient, $\pi=\kappa\varrho$, the properties of the symbols $\Res_S[\,\cdot\,/\Phi_{(\nu)}]$ and the
derivative relation as hypotheses, is checked by the Lean kernel (\texttt{PointwiseFormula.lean}, with Mathlib;
eleven theorems; the spanning statement is a hypothesis there). Referee~24 checked the spanning statement by
computer in $21$ degree types (re-run by the coordinating thread), and the derivative rule with a transcendental
parameter.
\end{remark}

\section{Open points, and what a specialist should check first}\label{sec:open}

\begin{enumerate}
\item \emph{Nothing here has been read by a human.} The order in which a specialist should check: (i) the cited
analytic input, on which everything rests: (G) with the scheme structure of the Hodge locus of a non-primitive
class and with its exact second-order coefficient, (A) and (SI) for complete intersections $Y_i$ that may be
singular, non-reduced or reducible; then Lemma~\ref{lem:PF}, in particular its new parts (3), (4), which replace
the cup-product formula (CG),
and hence Lemma~\ref{lem:T}(i) ($TV_i=(\Phi^i)_d$ and smoothness of $V_i$); (ii) Lemma~\ref{lem:mu} (found
necessary by referee~10); (iii) the key lemma, Step~3 of \S\ref{sec:G}; (iv) Lemma~\ref{lem:NW} and the
additivity \eqref{eq:link} of cycles under linkage, which were written by referees (they were part of the
versions read cold by later referees, and were re-derived by referee~22);
(v) Theorem~\ref{thm:FA}, whose proof is one page and would remove a hypothesis from \cite{P2}; (vi) the
range-free parts of Theorem~2, in particular (3) (which, as digested, is stronger than the published
statements) and (4) for cubics (which, as digested, is a published conjecture of Movasati, and which a literal
reading of one published theorem would contradict for $k\ge5$ at split hypersurfaces): each rests on ten lines
and on $TV_i=(\Phi^i)_d$, and should be tested against the literature first.
\item \emph{Novelty.} \S\ref{sec:lit} is a finite, machine-read search. In particular someone should read
\cite[\S6]{Kl} and say whether the liaison identity of Theorem~1 is implicit there beyond the linear case,
whether his cycles are the $W_t$ of \S\ref{sec:W}, and whether ``the Hodge locus of a linked determinantal
scheme is its deformation locus'' is recorded elsewhere.
\item \emph{The formula for $A$} (Theorem~\ref{thm:FA}, Corollary~\ref{cor:Anz}) rests on the period formula with a
constant independent of $X$, i.e. on Lemma~\ref{lem:PF}(4) (up to version~19: on (CG); alternatively on Villaflor
Loyola's theorem as extracted by a machine), and on the key lemma; Theorem~\ref{thm:FA} was read by AI referee~15 and re-derived by AI referees~16 and~22; its extension
Theorem~\ref{thm:FA2} was read by AI referees~17, 18 and~22, and Corollary~\ref{cor:surf} by referees~18 and~22.
\item \emph{Existence.} For some degree types no smooth $X$ exists; for the types of
Tables~\ref{tab:tau}--\ref{tab:types} existence is by random instances modulo primes.
\item \emph{$\et\ge2$, $\lambda\ne\lz$.} Theorem~\ref{thm:FA2} determines $\NL(\delta_\lambda)$ when the quadrics
have no common zero, which needs $\dim K\ge\et$. When they have common zeros (always if $\dim K<\et$, and at
special points such as the Fermat cubic with planes of the same matching) the structure is not determined here.
\item \emph{Outside the range} Theorems~1 and~2 hold, and Theorem~\ref{thm:FA3} gives the second-order law on
$E'$ against the forms of $K$. Not determined here: the second-order terms against $K_\lambda\smallsetminus K$
when (S) fails; and the structure when the quadrics on $\PP(E')$ have common zeros. The
formula of Theorem~\ref{thm:FA3}(2) rests on the same period formula as Theorem~\ref{thm:FA}.
\item \emph{Theorem~\ref{thm:lin}} rests, for its first-order half, on \S\S\ref{sec:facts}--\ref{sec:thm2} and on
the Lefschetz property of monomial complete intersections (\cite{St,Wa}, not opened; checked directly for
$d\le10$, $k\le9$); for its non-reduced half,
on Theorem~\ref{thm:FA3}, on Lemma~\ref{lem:PF}(4), and on exact computations with restricted data (reproduced by two referees,
with instances checkable by hand).
\item \emph{The Hodge conjecture.} Nothing in this note bears on it: all classes involved are algebraic.
\end{enumerate}

\begin{thebibliography}{99}\small
\bibitem{BH} W. Bruns, J. Herzog, \emph{Cohen--Macaulay rings}, Cambridge Univ. Press (not opened; Lemma~1.2.4 and
Prop.~1.2.13 are cited on the word of referees~11 and~12).
\bibitem{CG} J. Carlson, P. Griffiths, \emph{Infinitesimal variations of Hodge structure and the global Torelli
problem}, in: Journ\'ees de g\'eom\'etrie alg\'ebrique d'Angers 1979, 51--76 (not opened).
\bibitem{CPV} A. Corso, C. Polini, W. Vasconcelos, \emph{Links of prime ideals}, arXiv:math/0210034 (Theorem~2.5
as extracted by one agent; title as on the arXiv listing).
\bibitem{Dan1} A. Dan, \emph{Non-reduced components of the Noether--Lefschetz locus}, arXiv:1407.8491 (abstract
opened; partly machine-read).
\bibitem{Dan2} A. Dan, \emph{On a conjecture of Harris}, arXiv:2204.08079; Commun. Contemp. Math. (abstract
opened; machine-read).
\bibitem{DFV1} J. Duque Franco, R. Villaflor Loyola, \emph{On fake linear cycles inside Fermat varieties},
arXiv:2112.14818; Algebra Number Theory 17 (2023) (abstract only).
\bibitem{Fu} W. Fulton, \emph{Intersection theory}, Springer (not opened).
\bibitem{Gr} P. Griffiths, \emph{On the periods of certain rational integrals I, II}, Ann. of Math. 90 (1969)
(not opened).
\bibitem{Kl} R. Kloosterman, \emph{On a conjecture on Hodge loci of linear combinations of linear subvarieties},
arXiv:2312.12363v2; Rend. Circ. Mat. Palermo (2) 74 (2025), art.~187; Theorems~1.2, 1.3, 1.5, 4.9, 5.3,
Example~3.17, Prop.~4.1, 6.1, 6.4, Cor.~6.2, 6.5 (machine-read by three agents); Theorem~1.4, Definition~5.1 and
Lemma~2.11 (machine-read by one or two agents).
\bibitem{Kl2} R. Kloosterman, \emph{Hodge loci associated with linear subspaces intersecting in codimension
one}, arXiv:2401.10775; Math. Nachr. 298 (2025), 1220--1229; Theorem~1.1, Lemma~2.12, Theorem~2.13,
Propositions~3.1, 4.2, Theorem~4.8, Remark~4.9 (arXiv version machine-read by one agent; the journal version
was not read).
\bibitem{Mac} C. Maclean, \emph{A second-order invariant of the Noether--Lefschetz locus and two applications},
Asian J. Math. 9(3) (2005), 373--399 (not opened; her thesis was partly machine-read by one agent).
\bibitem{Mo} H. Movasati, \emph{On a Hodge locus}, with an appendix by R. Villaflor Loyola, arXiv:2211.11405
(machine-read).
\bibitem{Mo19} H. Movasati, \emph{Hodge cycles for cubic hypersurfaces}, arXiv:1902.00831 (machine-read).
\bibitem{Mo26} H. Movasati, \emph{On a counterexample to a conjecture of J. Harris for octic surfaces},
arXiv:2606.26944 (25 June 2026) (abstract opened; body machine-read).
\bibitem{No} D. G. Northcott, \emph{A homological investigation of a certain residual ideal}, Math. Ann. 150
(1963), 99--110 (not opened; data from \cite{CPV} as extracted).
\bibitem{P1} \emph{Hodge loci of sums of two linear cycles: Movasati's conjecture for cubic sixfolds and
eightfolds}, AI-produced manuscript, version 6, 21 November 2026.
\bibitem{P2} \emph{Noether--Lefschetz loci of sums of two disjoint complete-intersection curves}, AI-produced
manuscript, version 4, 21 November 2026.
\bibitem{PS} C. Peskine, L. Szpiro, \emph{Liaison des vari\'et\'es alg\'ebriques I}, Invent. Math. 26 (1974)
(not opened).
\bibitem{SS} A.-M. Simon, J. R. Strooker, \emph{Complete intersections of dimension zero: variations on a theme
of Wiebe}, arXiv:math/0703880 (Prop.~1.2 as extracted).
\bibitem{St} R. P. Stanley, \emph{Weyl groups, the hard Lefschetz theorem, and the Sperner property}, SIAM J.
Algebraic Discrete Methods 1 (1980), 168--184 (cited from memory by referee~21; not opened).
\bibitem{VL} R. Villaflor Loyola, \emph{Periods of complete intersection algebraic cycles}, Manuscripta Math.
167 (2022), 765--792; arXiv:1812.03964 (abstract and machine extraction; the period formula is Theorem~2 in the
arXiv version~5 and Theorem~1.2 in another rendering; volume and pages not confirmed by us).
\bibitem{Wa} J. Watanabe, \emph{The Dilworth number of Artinian rings and finite posets with rank function},
Adv. Stud. Pure Math. 11 (1987), 303--312 (cited from memory by referee~21; not opened).
\bibitem{Wi} H. Wiebe, \emph{\"Uber homologische Invarianten lokaler Ringe}, Math. Ann. 179 (1969), 257--274
(not opened; data from \cite{SS} as extracted).
\end{thebibliography}
\end{document}
